% Luc Destombe --- licence CC BY-NC-SA 4.0
% https://creativecommons.org/licenses/by-nc-sa/4.0/deed.fr
% Fichier autonome : compiler avec pdflatex, deux fois (signets, nombre de pages).
\documentclass[11pt,a4paper]{article}
\makeatletter
\newif\iflycee@unecolonne
\lycee@unecolonnetrue
\newif\iflycee@palatino
\lycee@palatinotrue

\RequirePackage[utf8]{inputenc}
\RequirePackage[T1]{fontenc}
\RequirePackage[french]{babel}
\RequirePackage{etoolbox}

\newcommand{\ShorthandsOff}{\@ifundefined{shorthandoff}{}{\shorthandoff{;:!?}}}
\newcommand{\ShorthandsOn}{\@ifundefined{shorthandon}{}{\shorthandon{;:!?}}}
\ShorthandsOff
\AtBeginDocument{\ShorthandsOn}
\AtBeginEnvironment{tikzpicture}{\ShorthandsOff}

\RequirePackage{geometry}
\RequirePackage{fancyhdr}
\RequirePackage{lastpage}
\RequirePackage{multicol}
\RequirePackage{array}
\RequirePackage{enumitem}
\RequirePackage{xcolor}
\RequirePackage{needspace}

\RequirePackage{amsmath,amssymb}
\RequirePackage{mathpazo}
\DeclareMathSizes{10.95}{10.95}{8}{5.5}
\begingroup\catcode`\_=\active \catcode`\^=\active
\gdef_#1{\sb{\scriptscriptstyle #1}}
\gdef^#1{\sp{\scriptscriptstyle\lycee@moinsepais #1}}
\endgroup
\newcommand\lycee@moins{\mathbin{%
  \setbox\z@\hbox{$\m@th\scriptscriptstyle\lycee@moinsorig$}%
  \dimen@=.55\wd\z@ \advance\dimen@-.5pt
  \hbox to.75\wd\z@{\hss$\m@th\scriptscriptstyle\vcenter{\hbox{%
    \tikz\draw[line width=.5pt,line cap=round](0,0)--(\the\dimen@,0);}}$\hss}}}
\begingroup\lccode`\~=`\- \lowercase{\endgroup\let~}\lycee@moins
\newcommand\lycee@moinsepais{\mathcode`\-="8000 }
\AtBeginDocument{\mathchardef\lycee@moinsorig=\mathcode`\-\relax}
\AtBeginDocument{\catcode`\_=12 \mathcode`\_="8000
  \catcode`\^=12 \mathcode`\^="8000 }
\newcommand\lycee@exposants{%
  \fontdimen13\textfont2=.48\fontdimen6\textfont2
  \fontdimen14\textfont2=.48\fontdimen6\textfont2
  \fontdimen15\textfont2=.48\fontdimen6\textfont2
  \fontdimen13\scriptfont2=.48\fontdimen6\scriptfont2
  \fontdimen14\scriptfont2=.48\fontdimen6\scriptfont2
  \fontdimen15\scriptfont2=.48\fontdimen6\scriptfont2}
\every@math@size\expandafter{\the\every@math@size\lycee@exposants}
\newlength{\retraitcalculs}
\setlength{\retraitcalculs}{2em}

\RequirePackage{tikz}
\usetikzlibrary{arrows.meta,calc,babel,shapes.misc}
\RequirePackage[most]{tcolorbox}
\tcbuselibrary{skins,breakable}

\definecolor{coulDef}{HTML}{1F4E79}
\definecolor{coulProp}{HTML}{7030A0}
\definecolor{coulRem}{HTML}{C55A11}
\definecolor{coulEx}{HTML}{2E7D32}
\definecolor{coulExo}{HTML}{B03A2E}
\definecolor{coulPart}{HTML}{1F4E79}
\definecolor{coulMeth}{HTML}{1B6E4F}
\definecolor{coulCourbe}{HTML}{0B5ED7}
\definecolor{coulCourbeB}{HTML}{C00000}
\definecolor{coulCourbeC}{HTML}{2E7D32}
\definecolor{coulGrille}{HTML}{8C8C8C}
\definecolor{coulRep}{HTML}{1B6E4F}
\definecolor{coulBar}{HTML}{2E7D32}

\newcommand{\R}{\mathbb{R}}
\newcommand{\Rs}{\mathbb{R}^{*}}
\renewcommand{\le}{\leqslant}
\renewcommand{\ge}{\geqslant}
\renewcommand{\approx}{\simeq}
\newcommand{\vect}[1]{\overrightarrow{#1}}
\newcommand{\norme}[1]{\left\lVert #1 \right\rVert}
\newcommand{\intervalle}[2]{\left[#1\,;#2\right]}
\RequirePackage{cancel}
\renewcommand{\CancelColor}{\color{coulExo}}
\newcommand{\fois}[1]{\times{\color{coulExo}#1}}
\newcommand{\barre}[1]{\cancel{#1}}

\tikzset{grille/.style={coulGrille,line width=0.4pt}}
\newcommand{\quadrillage}[4]{\draw[grille,step=1] (#1,#2) grid (#3,#4);}
\tikzset{
  croix/.style={cross out,draw,line width=0.6pt,inner sep=0pt,minimum size=4pt},
  croixdroite/.style={croix,rotate=45,line width=0.8pt},
}
\newcommand{\pt}[3]{\path (#1) node[croix]{} node[#2,font=\small,inner sep=2.5pt]{$#3$};}

\newcommand{\lycee@repere}[4]{%
  \draw[grille,step=1] (#1,#2) grid (#3,#4);
  \draw[-{Stealth[length=2mm]},thick] (#1,0)--({#3+0.4},0) node[below left=-1pt]{$x$};
  \draw[-{Stealth[length=2mm]},thick] (0,#2)--(0,{#4+0.4}) node[below left=-1pt]{$y$};
  \node[below left,font=\scriptsize,inner sep=1.5pt] at (0,0) {$0$};}
\newcommand{\lycee@gradx}[1]{\foreach \k/\l in {#1}{%
  \draw (\k,0.07)--(\k,-0.07) node[below,font=\scriptsize,inner sep=1.5pt]{$\l$};}}
\newcommand{\lycee@grady}[1]{\foreach \k/\l in {#1}{%
  \draw (0.07,\k)--(-0.07,\k) node[left,font=\scriptsize,inner sep=1.5pt]{$\l$};}}
\newcommand{\lycee@intff}[2]{\left[#1\,;#2\right]}
\newcommand{\lycee@intfo}[2]{\left[#1\,;#2\right[}
\newcommand{\lycee@intof}[2]{\left]#1\,;#2\right]}
\newcommand{\lycee@intoo}[2]{\left]#1\,;#2\right[}
\newcommand{\lycee@ens}[1]{\left\{#1\right\}}
\AtBeginDocument{%
  \providecommand{\repere}{\lycee@repere}%
  \providecommand{\gradx}{\lycee@gradx}%
  \providecommand{\grady}{\lycee@grady}%
  \providecommand{\Cf}{\mathcal{C}\sb{\scriptscriptstyle f}}%
  \providecommand{\Cg}{\mathcal{C}\sb{\scriptscriptstyle g}}%
  \providecommand{\intff}{\lycee@intff}%
  \providecommand{\intfo}{\lycee@intfo}%
  \providecommand{\intof}{\lycee@intof}%
  \providecommand{\intoo}{\lycee@intoo}%
  \providecommand{\ens}{\lycee@ens}}

\newlist{questions}{enumerate}{3}
\setlist[questions,1]{label=\arabic*\textdegree),leftmargin=*,itemsep=2pt,topsep=3pt}
\setlist[questions,2]{label=\alph*),leftmargin=*,itemsep=1pt,topsep=2pt}
\setlist[questions,3]{label=\roman*),leftmargin=*,itemsep=1pt,topsep=1pt}
\setlist[itemize]{label=\textbullet,leftmargin=*,itemsep=2pt,topsep=3pt}

\newcommand{\besoinplace}[1]{\par\penalty\@M\begingroup
  \dimen@=#1\relax\dimen@ii=\pagegoal\advance\dimen@ii-\pagetotal
  \ifdim\dimen@>\dimen@ii\ifdim\dimen@ii>\z@\vfil\fi\break\fi\endgroup}

\newcommand{\lycee@pied}{\thepage/\pageref{LastPage}}
\newcommand{\entete}[2]{%
  \pagestyle{fancy}\fancyhf{}%
  \fancyhead[L]{\footnotesize\color{coulPart}\textbf{#1}}%
  \fancyhead[R]{\footnotesize\color{coulPart}\textbf{#2}}%
  \fancyfoot[C]{\footnotesize\lycee@pied}%
  \renewcommand{\headrulewidth}{0.4pt}%
  \renewcommand{\headrule}{\hbox to\headwidth{%
    \color{coulPart!40}\leaders\hrule height \headrulewidth\hfill}}}

\RequirePackage{ccicons}
\newcommand{\lycee@licence}{{\scriptsize\color{gray}%
  \copyright\ Luc Destombe --- \ccbyncsa\ CC BY-NC-SA 4.0}}
\newif\iflycee@licence \lycee@licencetrue
\newcommand{\sanslicence}{\lycee@licencefalse}
\AtBeginDocument{\iflycee@licence\AddToHookNext{shipout/foreground}{%
  \put(0,0){\rlap{\kern\dimexpr1in+\hoffset+\oddsidemargin+\textwidth\relax
    \llap{\raisebox{-\dimexpr1in+\voffset+\topmargin+\headheight+\headsep
      +\textheight+\footskip\relax}{\lycee@licence}}}}}\fi}

\newcommand{\formulesserrees}{%
  \setlength{\abovedisplayskip}{3pt}\setlength{\belowdisplayskip}{3pt}%
  \setlength{\abovedisplayshortskip}{2pt}\setlength{\belowdisplayshortskip}{3pt}}

\setlength{\parindent}{0pt}

\RequirePackage[implicit=false,hidelinks,pdfencoding=auto,psdextra,bookmarksopen,
  bookmarksopenlevel=1,pdfcreator={},pdfproducer={}]{hyperref}
\RequirePackage{bookmark}
\hypersetup{pdfauthor={Luc Destombe},pdfkeywords={CC BY-NC-SA 4.0}}
\newcounter{lycee@signet}
\newcommand{\signet}[3][]{\stepcounter{lycee@signet}%
  \edef\lycee@dest{\if\relax\detokenize{#1}\relax
    signet.\arabic{lycee@signet}\else#1\fi}%
  \hypertarget{\lycee@dest}{}\bookmark[level=#2,dest=\lycee@dest]{#3}}
\pdfstringdefDisableCommands{%
  \def\R{R}\def\vect#1{#1}\def\Cf{Cf}\def\Cg{Cg}\def\fois#1{×#1}%
  \def\barre#1{#1}\def\intervalle#1#2{[#1;#2]}\def\intff#1#2{[#1;#2]}%
  \def\textsuperscript#1{#1}\def\dfrac#1#2{#1/#2}\def\frac#1#2{#1/#2}%
  \def\vec#1{#1⃗}}%
\RequirePackage[official]{eurosym}

\tcbset{exercice corrige/.style={
  enhanced,breakable,before upper=\raggedright,
  colback=white,colframe=coulExo,
  boxrule=0.8pt,arc=2pt,
  left=6pt,right=6pt,top=8pt,bottom=5pt,
  before skip=9pt,after skip=4pt,
  coltitle=white,fonttitle=\bfseries\small,
  attach boxed title to top left={xshift=8pt,yshift=-\tcboxedtitleheight/2},
  boxed title style={colback=coulExo,arc=2pt,boxrule=0pt}}}
\newcounter{exo}
\newtcolorbox{exercice}[1][]{exercice corrige,
  phantom={\signet[exercice-\arabic{exo}]{2}{Exercice \arabic{exo}}},
  title={Exercice \theexo\ifx\relax#1\relax\else{} --- #1\fi}}
\newenvironment{exo}[1][\relax]{\refstepcounter{exo}\begin{exercice}[#1]}{\end{exercice}}

\geometry{top=1.8cm,bottom=1cm,left=1cm,right=1cm,headheight=14pt,footskip=0.6cm}
\renewcommand{\lycee@pied}{\thepage}

\renewcommand{\theexo}{\ifnum\value{exo}<10 0\fi\arabic{exo}}
\tcbset{exercice corrige/.style={
  enhanced,breakable,before upper=\raggedright,
  colback=white,colframe=coulExo,
  boxrule=0.9pt,arc=2pt,
  left=8pt,right=8pt,top=8pt,bottom=6pt,
  before skip=8pt,after skip=6pt,
  coltitle=white,fonttitle=\bfseries,
  attach boxed title to top left={xshift=8pt,yshift=-\tcboxedtitleheight/2},
  boxed title style={colback=coulExo,arc=2pt,boxrule=0pt}}}

\newcommand{\entetefiche}[2]{%
  \begin{center}
    \begin{tcolorbox}[enhanced,width=0.92\linewidth,colback=white,colframe=coulPart,
        boxrule=1.4pt,arc=3pt,halign=center,top=6pt,bottom=6pt]
      {\large\bfseries\color{coulPart} CORRIGÉ --- FICHE D'EXERCICES}\\[3pt]
      {\normalsize\bfseries #1}\\[2pt]
      {\normalsize #2}
    \end{tcolorbox}
  \end{center}}

\newcounter{partie}
\newcommand{\partie}[1]{%
  \stepcounter{partie}%
  \besoinplace{8\baselineskip}\vspace{4pt}%
  \begin{tcolorbox}[enhanced,colback=coulPart!8,colframe=coulPart,boxrule=0pt,
      leftrule=3pt,arc=0pt,left=6pt,right=4pt,top=3pt,bottom=3pt,
      before skip=6pt,after skip=2pt]
    \bfseries\color{coulPart}\Roman{partie} --- #1\signet{1}{\Roman{partie} --- #1}
  \end{tcolorbox}}
\newcommand{\partiesansnum}[1]{%
  \besoinplace{8\baselineskip}\vspace{4pt}%
  \begin{tcolorbox}[enhanced,colback=coulPart!8,colframe=coulPart,boxrule=0pt,
      leftrule=3pt,arc=0pt,left=6pt,right=4pt,top=3pt,bottom=3pt,
      before skip=6pt,after skip=2pt]
    \bfseries\color{coulPart}#1\signet{1}{#1}
  \end{tcolorbox}}

\setlist[questions,1]{label=\arabic*\textdegree),leftmargin=*,itemsep=3pt,topsep=3pt}
\setlist[questions,2]{label=\alph*),leftmargin=*,itemsep=2pt,topsep=2pt}
\setlist[questions,3]{label=\roman*),leftmargin=*,itemsep=1pt,topsep=1pt}
\newlist{expr}{enumerate}{1}
\setlist[expr]{label={\Alph*\ $=$},leftmargin=*,itemsep=2pt,topsep=3pt}

\newcommand{\res}[1]{{\color{coulRep}#1}}
\newcommand{\rem}[1]{\par\smallskip{\small\itshape\color{black!60}$\rightarrow$ #1}}
\newenvironment{calculs}{\par\smallskip\noindent$\displaystyle\begin{aligned}[t]}%
  {\end{aligned}$\par\smallskip}
\newcommand{\justif}[1]{\quad\text{\small\itshape\color{black!60}#1}}
\newcommand{\cas}[1]{\par\smallskip\noindent\textbf{\color{coulExo}#1}\nobreak}
\newcommand{\explique}[2]{\par\smallskip\noindent
  \begin{minipage}[t]{0.6\linewidth}\raggedright #1\end{minipage}\hfill
  \begin{minipage}[t]{0.37\linewidth}\raggedright\leavevmode
    {\small\itshape\color{black!60}#2}\end{minipage}\par}

\setlength{\parskip}{2pt}

\makeatother
\usepackage{tkz-tab}
\entete{Chapitre 3 --- Calcul littéral et second degré --- corrigé}{1\textsuperscript{re} mathématiques spécifiques}
\setlength{\parindent}{0pt}
\setlength{\jot}{5pt}

\let\casligne\cas
\renewcommand{\cas}[1]{\casligne{#1}\par\nopagebreak}

\newcommand{\varmin}[3]{\par\smallskip{\centering
  \begin{tikzpicture}[scale=0.85]
  \tkzTabInit[lgt=1.6,espcl=2.2]{$x$/1,$#1$/1.6}{,$#2$,}
  \tkzTabVar{+/ ,-/$#3$,+/ }
  \end{tikzpicture}\par}\smallskip}
\newcommand{\varmax}[3]{\par\smallskip{\centering
  \begin{tikzpicture}[scale=0.85]
  \tkzTabInit[lgt=1.6,espcl=2.2]{$x$/1,$#1$/1.6}{,$#2$,}
  \tkzTabVar{-/ ,+/$#3$,-/ }
  \end{tikzpicture}\par}\smallskip}
\newcommand{\varint}[4][x]{\par\smallskip{\centering
  \begin{tikzpicture}[scale=0.85]
  \tkzTabInit[lgt=1.6,espcl=2.2]{$#1$/1,$#2$/1.6}{#3}
  \tkzTabVar{#4}
  \end{tikzpicture}\par}\smallskip}

\begin{document}

\entetefiche{Chapitre 3 --- Calcul littéral et second degré}{1\textsuperscript{re}
mathématiques spécifiques}

\partie{Développer}

\begin{exo}[Distributivité simple]
\explique{Dans $k(a+b)$, le nombre devant la parenthèse multiplie \emph{chaque} terme de
la parenthèse : $k(a+b)=ka+kb$.}{méthode 1 du cours}
\begin{calculs}
A &= 5(x+3) &&\justif{on recopie}\\
A &= 5\times x+5\times 3 &&\justif{$5$ multiplie $x$, puis $3$}\\
A &= \res{5x+15} &&
\end{calculs}
\begin{calculs}
B &= 4(2x-7) &&\justif{on recopie}\\
B &= 4\times 2x-4\times 7 &&\justif{$4$ multiplie $2x$, puis $7$}\\
B &= \res{8x-28} &&\justif{$4\times 2x=8x$}
\end{calculs}
\begin{calculs}
C &= -3(x+6) &&\justif{on recopie}\\
C &= (-3)\times x+(-3)\times 6 &&\justif{c'est $-3$, avec son signe, qui multiplie}\\
C &= \res{-3x-18} &&\justif{$(-3)\times 6=-18$}
\end{calculs}
\begin{calculs}
D &= -2(4-5x) &&\justif{on recopie}\\
D &= (-2)\times 4+(-2)\times(-5x) &&\justif{le terme $-5x$ garde son signe}\\
D &= -8+10x &&\justif{moins par moins donne plus : $(-2)\times(-5x)=+10x$}\\
D &= \res{10x-8} &&\justif{on range : le terme en $x$ d'abord}
\end{calculs}
\begin{calculs}
E &= 7(3-x) &&\justif{on recopie}\\
E &= 7\times 3+7\times(-x) &&\\
E &= \res{21-7x} &&
\end{calculs}
\begin{calculs}
F &= -(x-9) &&\justif{le signe $-$ devant la parenthèse, c'est $-1$}\\
F &= (-1)\times x+(-1)\times(-9) &&\justif{$-1$ multiplie chaque terme}\\
F &= \res{-x+9} &&\justif{un $-$ devant une parenthèse change tous les signes}
\end{calculs}
\rem{Erreur fréquente : ne multiplier que le premier terme, $5(x+3)=5x+3$. Le nombre
devant la parenthèse multiplie les \emph{deux} termes.}
\rem{Avec un nombre négatif devant, on écrit les produits avec des parenthèses, comme
$(-2)\times(-5x)$ : on voit alors la règle des signes.}
\end{exo}

\begin{exo}[Un facteur en $x$]
\explique{Même méthode : le facteur devant la parenthèse, ici un terme en $x$, multiplie
chaque terme.}{$x\times x=x^{2}$}
\begin{calculs}
A &= x(x+8) &&\justif{on recopie}\\
A &= x\times x+x\times 8 &&\justif{$x$ multiplie chaque terme}\\
A &= \res{x^{2}+8x} &&\justif{$x\times x=x^{2}$}
\end{calculs}
\begin{calculs}
B &= 2x(3x-1) &&\justif{on recopie}\\
B &= 2x\times 3x-2x\times 1 &&\justif{$2x$ multiplie chaque terme}\\
B &= \res{6x^{2}-2x} &&\justif{$2x\times 3x=2\times 3\times x\times x=6x^{2}$}
\end{calculs}
\begin{calculs}
C &= -x(x-4) &&\justif{on recopie}\\
C &= (-x)\times x+(-x)\times(-4) &&\justif{$-x$ multiplie chaque terme}\\
C &= \res{-x^{2}+4x} &&\justif{$(-x)\times(-4)=+4x$}
\end{calculs}
\begin{calculs}
D &= -4x(2x+5) &&\justif{on recopie}\\
D &= (-4x)\times 2x+(-4x)\times 5 &&\justif{$-4x$ multiplie chaque terme}\\
D &= \res{-8x^{2}-20x} &&\justif{$(-4)\times 2=-8$ et $x\times x=x^{2}$}
\end{calculs}
\begin{calculs}
E &= 3x(5-2x) &&\justif{on recopie}\\
E &= 3x\times 5+3x\times(-2x) &&\justif{$3x$ multiplie chaque terme}\\
E &= \res{15x-6x^{2}} &&\justif{$3x\times(-2x)=-6x^{2}$}
\end{calculs}
\rem{On multiplie les nombres entre eux, puis les $x$ entre eux : $2x\times 3x=6x^{2}$,
et non $6x$ ni $5x^{2}$.}
\end{exo}

\begin{exo}[Double distributivité]
\explique{Dans $(a+b)(c+d)$, chaque terme de la première parenthèse multiplie chaque
terme de la seconde : quatre produits, puis on réduit.}{méthode 1 du cours}
\begin{calculs}
A &= (x+2)(x+5) &&\justif{on recopie}\\
A &= x\times x+x\times 5+2\times x+2\times 5 &&\justif{quatre produits}\\
A &= x^{2}+5x+2x+10 &&\\
A &= \res{x^{2}+7x+10} &&\justif{on réduit : $5x+2x=7x$}
\end{calculs}
\begin{calculs}
B &= (x-3)(x+4) &&\justif{on recopie}\\
B &= x\times x+x\times 4+(-3)\times x+(-3)\times 4 &&\justif{quatre produits ; $-3$ garde
son signe}\\
B &= x^{2}+4x-3x-12 &&\\
B &= \res{x^{2}+x-12} &&\justif{$4x-3x=1x$, qu'on écrit $x$}
\end{calculs}
\begin{calculs}
C &= (x-6)(x-1) &&\justif{on recopie}\\
C &= x\times x+x\times(-1)+(-6)\times x+(-6)\times(-1) &&\justif{quatre produits}\\
C &= x^{2}-x-6x+6 &&\justif{$(-6)\times(-1)=+6$}\\
C &= \res{x^{2}-7x+6} &&\justif{$-x-6x=-7x$}
\end{calculs}
\begin{calculs}
D &= (2x+1)(x+3) &&\justif{on recopie}\\
D &= 2x\times x+2x\times 3+1\times x+1\times 3 &&\justif{quatre produits}\\
D &= 2x^{2}+6x+x+3 &&\\
D &= \res{2x^{2}+7x+3} &&\justif{$6x+x=7x$}
\end{calculs}
\begin{calculs}
E &= (3x-2)(x+2) &&\justif{on recopie}\\
E &= 3x\times x+3x\times 2+(-2)\times x+(-2)\times 2 &&\justif{quatre produits}\\
E &= 3x^{2}+6x-2x-4 &&\\
E &= \res{3x^{2}+4x-4} &&\justif{$6x-2x=4x$}
\end{calculs}
\begin{calculs}
F &= (x+5)(4-x) &&\justif{on recopie}\\
F &= x\times 4+x\times(-x)+5\times 4+5\times(-x) &&\justif{quatre produits}\\
F &= 4x-x^{2}+20-5x &&\justif{$x\times(-x)=-x^{2}$}\\
F &= \res{-x^{2}-x+20} &&\justif{$4x-5x=-x$ ; on range : $x^{2}$, puis $x$, puis le
nombre}
\end{calculs}
\rem{Pour n'oublier aucun produit, on compte : il en faut quatre. On peut aussi relier
par des flèches chaque terme de la première parenthèse à chaque terme de la seconde.}
\rem{Contrôle rapide : on remplace $x$ par $1$ des deux côtés. Pour $A$ :
$(1+2)\times(1+5)=18$ et $1+7+10=18$.}
\end{exo}

\begin{exo}[Vrai ou faux ?]
\explique{Pour savoir si une égalité est vraie pour tout $x$, on développe le membre de
gauche et on compare avec celui de droite.}{on part du côté avec des parenthèses}
\cas{a)}
\begin{calculs}
3(x+4) &= 3\times x+3\times 4 &&\justif{$3$ multiplie chaque terme}\\
3(x+4) &= 3x+12 &&
\end{calculs}
\explique{\res{Faux} : on trouve $3x+12$, pas $3x+4$. L'égalité correcte est
$3(x+4)=\res{3x+12}$.}{le $4$ a été oublié dans la multiplication}
\cas{b)}
\begin{calculs}
-2(x-5) &= (-2)\times x+(-2)\times(-5) &&\justif{$-2$ multiplie chaque terme}\\
-2(x-5) &= -2x+10 &&\justif{$(-2)\times(-5)=+10$}
\end{calculs}
\explique{\res{Vrai}.}{}
\cas{c)}
\begin{calculs}
x(x+1) &= x\times x+x\times 1 &&\justif{$x$ multiplie chaque terme}\\
x(x+1) &= x^{2}+x &&
\end{calculs}
\explique{\res{Vrai}.}{}
\cas{d)}
\begin{calculs}
(x+1)(x+2) &= x\times x+x\times 2+1\times x+1\times 2 &&\justif{quatre produits}\\
(x+1)(x+2) &= x^{2}+2x+x+2 &&\\
(x+1)(x+2) &= x^{2}+3x+2 &&\justif{$2x+x=3x$}
\end{calculs}
\explique{\res{Faux} : l'égalité correcte est $(x+1)(x+2)=\res{x^{2}+3x+2}$.}{il manque
les deux produits \og croisés\fg{} $x\times 2$ et $1\times x$}
\explique{Un seul exemple suffit pour montrer qu'elle est fausse : pour $x=1$,
$(1+1)\times(1+2)=6$ mais $1^{2}+2=3$.}{tester une valeur}
\cas{e)}
\begin{calculs}
-(3-x) &= (-1)\times 3+(-1)\times(-x) &&\justif{le $-$ devant, c'est $-1$}\\
-(3-x) &= -3+x &&\justif{$(-1)\times(-x)=+x$}
\end{calculs}
\explique{\res{Faux} : l'égalité correcte est $-(3-x)=\res{-3+x}$.}{un $-$ devant la
parenthèse change \emph{tous} les signes}
\rem{Erreur fréquente (e) : changer le signe du premier terme seulement.}
\end{exo}

\partie{Identités remarquables}

\begin{exo}[Le carré d'une somme]
\explique{On reconnaît $(a+b)^{2}$ : un carré, avec un $+$ dans la parenthèse. On
repère $a$ et $b$, puis on applique $(a+b)^{2}=a^{2}+2ab+b^{2}$.}{méthode 2 du cours}
\begin{calculs}
A &= (x+4)^{2} &&\justif{$(a+b)^{2}$, avec $a=x$ et $b=4$}\\
A &= x^{2}+2\times x\times 4+4^{2} &&\justif{$a^{2}+2ab+b^{2}$}\\
A &= \res{x^{2}+8x+16} &&\justif{$2\times 4=8$ et $4^{2}=16$}
\end{calculs}
\begin{calculs}
B &= (x+1)^{2} &&\justif{$a=x$ et $b=1$}\\
B &= x^{2}+2\times x\times 1+1^{2} &&\justif{$a^{2}+2ab+b^{2}$}\\
B &= \res{x^{2}+2x+1} &&
\end{calculs}
\begin{calculs}
C &= (2x+3)^{2} &&\justif{$a=2x$ et $b=3$}\\
C &= (2x)^{2}+2\times 2x\times 3+3^{2} &&\justif{on garde les parenthèses autour de $2x$}\\
C &= \res{4x^{2}+12x+9} &&\justif{$(2x)^{2}=2x\times 2x=4x^{2}$ ; $2\times 2\times 3=12$}
\end{calculs}
\begin{calculs}
D &= (3x+5)^{2} &&\justif{$a=3x$ et $b=5$}\\
D &= (3x)^{2}+2\times 3x\times 5+5^{2} &&\justif{$a^{2}+2ab+b^{2}$}\\
D &= \res{9x^{2}+30x+25} &&\justif{$(3x)^{2}=9x^{2}$ ; $2\times 3\times 5=30$}
\end{calculs}
\rem{Erreur fréquente : oublier le double produit, $(x+4)^{2}=x^{2}+16$. Contrôle avec
$x=1$ : $(1+4)^{2}=25$, alors que $1+16=17$.}
\rem{Autre piège : $(2x)^{2}$ vaut $4x^{2}$, pas $2x^{2}$ ; le $2$ est lui aussi élevé
au carré.}
\end{exo}

\begin{exo}[Le carré d'une différence]
\explique{On reconnaît $(a-b)^{2}$ : un carré, avec un $-$ dans la parenthèse. On
applique $(a-b)^{2}=a^{2}-2ab+b^{2}$ ; seul le double produit est
négatif.}{méthode 2 du cours}
\begin{calculs}
A &= (x-3)^{2} &&\justif{$(a-b)^{2}$, avec $a=x$ et $b=3$}\\
A &= x^{2}-2\times x\times 3+3^{2} &&\justif{$a^{2}-2ab+b^{2}$}\\
A &= \res{x^{2}-6x+9} &&
\end{calculs}
\begin{calculs}
B &= (x-10)^{2} &&\justif{$a=x$ et $b=10$}\\
B &= x^{2}-2\times x\times 10+10^{2} &&\justif{$a^{2}-2ab+b^{2}$}\\
B &= \res{x^{2}-20x+100} &&
\end{calculs}
\begin{calculs}
C &= (2x-1)^{2} &&\justif{$a=2x$ et $b=1$}\\
C &= (2x)^{2}-2\times 2x\times 1+1^{2} &&\justif{$a^{2}-2ab+b^{2}$}\\
C &= \res{4x^{2}-4x+1} &&\justif{$(2x)^{2}=4x^{2}$}
\end{calculs}
\begin{calculs}
D &= (5x-2)^{2} &&\justif{$a=5x$ et $b=2$}\\
D &= (5x)^{2}-2\times 5x\times 2+2^{2} &&\justif{$a^{2}-2ab+b^{2}$}\\
D &= \res{25x^{2}-20x+4} &&\justif{$(5x)^{2}=25x^{2}$ ; $2\times 5\times 2=20$}
\end{calculs}
\rem{Le dernier terme $b^{2}$ est toujours positif : $(x-3)^{2}$ se termine par $+9$, pas
par $-9$.}
\end{exo}

\begin{exo}[Somme par différence]
\explique{On reconnaît $(a+b)(a-b)$ : deux parenthèses avec les mêmes termes, l'une avec
$+$, l'autre avec $-$. Le résultat est $a^{2}-b^{2}$, sans terme en $x$.}{méthode 2 du
cours}
\begin{calculs}
A &= (x+6)(x-6) &&\justif{$(a+b)(a-b)$, avec $a=x$ et $b=6$}\\
A &= x^{2}-6^{2} &&\justif{$a^{2}-b^{2}$}\\
A &= \res{x^{2}-36} &&
\end{calculs}
\begin{calculs}
B &= (x-1)(x+1) &&\justif{$a=x$ et $b=1$ ; l'ordre des parenthèses ne compte pas}\\
B &= x^{2}-1^{2} &&\justif{$a^{2}-b^{2}$}\\
B &= \res{x^{2}-1} &&
\end{calculs}
\begin{calculs}
C &= (3x+2)(3x-2) &&\justif{$a=3x$ et $b=2$}\\
C &= (3x)^{2}-2^{2} &&\justif{$a^{2}-b^{2}$}\\
C &= \res{9x^{2}-4} &&\justif{$(3x)^{2}=9x^{2}$}
\end{calculs}
\begin{calculs}
D &= (4x-5)(4x+5) &&\justif{$a=4x$ et $b=5$}\\
D &= (4x)^{2}-5^{2} &&\justif{$a^{2}-b^{2}$}\\
D &= \res{16x^{2}-25} &&
\end{calculs}
\rem{Pourquoi n'y a-t-il pas de terme en $x$ ? En faisant les quatre produits de $A$, on
trouve $+6x$ et $-6x$, qui s'annulent : $x^{2}-6x+6x-36=x^{2}-36$.}
\end{exo}

\begin{exo}[La bonne réponse]
\cas{1.}
\begin{calculs}
(x-7)^{2} &= x^{2}-2\times x\times 7+7^{2} &&\justif{$(a-b)^{2}$, avec $a=x$ et $b=7$}\\
(x-7)^{2} &= x^{2}-14x+49 &&
\end{calculs}
\explique{Réponse \res{b)}.}{a) et c) oublient le double produit ; d) oublie le $2$ du
double produit}
\cas{2.}
\begin{calculs}
(3x+1)^{2} &= (3x)^{2}+2\times 3x\times 1+1^{2} &&\justif{$(a+b)^{2}$, avec $a=3x$ et
$b=1$}\\
(3x+1)^{2} &= 9x^{2}+6x+1 &&
\end{calculs}
\explique{Réponse \res{c)}.}{a) oublie le double produit ; b) écrit $3x^{2}$ au lieu de
$9x^{2}$}
\cas{3.}
\explique{$x^{2}-16=x^{2}-4^{2}$ : c'est la forme $a^{2}-b^{2}$, avec $a=x$ et $b=4$.
Donc $x^{2}-16=(x-4)(x+4)$ : réponse \res{b)}.}{on lit l'identité de droite à gauche}
\explique{Contrôle : $(x-4)(x+4)=x^{2}+4x-4x-16=x^{2}-16$.}{a) donnerait
$x^{2}-8x+16$}
\cas{4.}
\begin{calculs}
(2x-5)(2x+5) &= (2x)^{2}-5^{2} &&\justif{$(a-b)(a+b)$, avec $a=2x$ et $b=5$}\\
(2x-5)(2x+5) &= 4x^{2}-25 &&
\end{calculs}
\explique{Réponse \res{a)}.}{b) oublie de mettre le $2$ au carré}
\rem{Dans un QCM, on calcule d'abord soi-même, puis on cherche son résultat parmi les
réponses : les mauvaises réponses sont les erreurs les plus fréquentes.}
\end{exo}

\partie{Équations du premier degré}

\begin{exo}[Premières équations]
\explique{On fait la même opération sur les deux membres de l'égalité : on regroupe les
$x$ d'un côté, les nombres de l'autre, puis on divise par le nombre devant
$x$.}{méthode 3 du cours}
\cas{a)}
\begin{calculs}
x+7 &= 12 &&\justif{on recopie}\\
x &= 12-7 &&\justif{on retranche $7$ aux deux membres}\\
x &= 5 &&
\end{calculs}
\explique{$\res{S=\ens{5}}$.}{contrôle : $5+7=12$}
\cas{b)}
\begin{calculs}
4x &= -20 &&\justif{on recopie}\\
x &= \dfrac{-20}{4} &&\justif{on divise les deux membres par $4$}\\
x &= -5 &&
\end{calculs}
\explique{$\res{S=\ens{-5}}$.}{contrôle : $4\times(-5)=-20$}
\cas{c)}
\begin{calculs}
3x-5 &= 10 &&\justif{on recopie}\\
3x &= 10+5 &&\justif{on ajoute $5$ aux deux membres}\\
3x &= 15 &&\\
x &= 5 &&\justif{on divise par $3$}
\end{calculs}
\explique{$\res{S=\ens{5}}$.}{contrôle : $3\times 5-5=10$}
\cas{d)}
\begin{calculs}
5x+2 &= 3x+10 &&\justif{on recopie}\\
5x-3x &= 10-2 &&\justif{les $x$ à gauche, les nombres à droite}\\
2x &= 8 &&\\
x &= 4 &&\justif{on divise par $2$}
\end{calculs}
\explique{$\res{S=\ens{4}}$.}{contrôle : $5\times 4+2=22$ et $3\times 4+10=22$}
\cas{e)}
\begin{calculs}
8x-1 &= 2x+11 &&\justif{on recopie}\\
8x-2x &= 11+1 &&\justif{$-2x$ passe à gauche, $-1$ passe à droite}\\
6x &= 12 &&\\
x &= 2 &&\justif{on divise par $6$}
\end{calculs}
\explique{$\res{S=\ens{2}}$.}{contrôle : $8\times 2-1=15$ et $2\times 2+11=15$}
\cas{f)}
\begin{calculs}
2x+9 &= 6x-3 &&\justif{on recopie}\\
9+3 &= 6x-2x &&\justif{les $x$ du côté où il y en a le plus}\\
12 &= 4x &&\\
x &= 3 &&\justif{on divise par $4$}
\end{calculs}
\explique{$\res{S=\ens{3}}$.}{contrôle : $2\times 3+9=15$ et $6\times 3-3=15$}
\rem{Un terme qui change de côté change de signe : c'est la même chose que retrancher
(ou ajouter) ce terme aux deux membres.}
\rem{En f), mettre les $x$ à droite évite un coefficient négatif devant $x$.}
\end{exo}

\begin{exo}[Avec des parenthèses]
\explique{On développe d'abord les parenthèses ; on retrouve alors une équation comme
celles de l'exercice 9.}{méthode 3 du cours}
\cas{a)}
\begin{calculs}
3(x+2) &= 21 &&\justif{on recopie}\\
3x+6 &= 21 &&\justif{on développe : $3\times x+3\times 2$}\\
3x &= 21-6 &&\justif{on retranche $6$}\\
3x &= 15 &&\\
x &= 5 &&\justif{on divise par $3$}
\end{calculs}
\explique{$\res{S=\ens{5}}$.}{contrôle : $3\times(5+2)=21$}
\cas{b)}
\begin{calculs}
2(x-4) &= x+1 &&\justif{on recopie}\\
2x-8 &= x+1 &&\justif{on développe}\\
2x-x &= 1+8 &&\justif{les $x$ à gauche, les nombres à droite}\\
x &= 9 &&
\end{calculs}
\explique{$\res{S=\ens{9}}$.}{contrôle : $2\times(9-4)=10$ et $9+1=10$}
\cas{c)}
\begin{calculs}
5(x+1) &= 2(x+7) &&\justif{on recopie}\\
5x+5 &= 2x+14 &&\justif{on développe les deux membres}\\
5x-2x &= 14-5 &&\\
3x &= 9 &&\\
x &= 3 &&\justif{on divise par $3$}
\end{calculs}
\explique{$\res{S=\ens{3}}$.}{contrôle : $5\times 4=20$ et $2\times 10=20$}
\cas{d)}
\begin{calculs}
-(x-6) &= 2x &&\justif{on recopie}\\
-x+6 &= 2x &&\justif{le $-$ devant la parenthèse change tous les signes}\\
6 &= 2x+x &&\justif{les $x$ à droite, où il y en a le plus}\\
6 &= 3x &&\\
x &= 2 &&\justif{on divise par $3$}
\end{calculs}
\explique{$\res{S=\ens{2}}$.}{contrôle : $-(2-6)=4$ et $2\times 2=4$}
\cas{e)}
\begin{calculs}
4(2x-1) &= 3(x+7) &&\justif{on recopie}\\
8x-4 &= 3x+21 &&\justif{on développe les deux membres}\\
8x-3x &= 21+4 &&\\
5x &= 25 &&\\
x &= 5 &&\justif{on divise par $5$}
\end{calculs}
\explique{$\res{S=\ens{5}}$.}{contrôle : $4\times 9=36$ et $3\times 12=36$}
\rem{Erreur fréquente : $3(x+2)=21$ devient $3x+2=21$ ; le $3$ multiplie aussi le $2$.}
\end{exo}

\begin{exo}[L'inconnue au dénominateur]
\explique{Quand $x$ est au dénominateur, on multiplie les deux membres par $x$ (qui n'est
pas nul) : $x$ remonte, et on se ramène à une équation du type $3x=12$.}{méthode 3 du
cours}
\cas{a)}
\begin{calculs}
\dfrac{12}{x} &= 3 &&\justif{on recopie}\\
12 &= 3x &&\justif{on multiplie les deux membres par $x$}\\
x &= 4 &&\justif{on divise par $3$}
\end{calculs}
\explique{$\res{S=\ens{4}}$.}{contrôle : $12\div 4=3$}
\cas{b)}
\begin{calculs}
\dfrac{30}{x} &= 6 &&\justif{on recopie}\\
30 &= 6x &&\justif{on multiplie par $x$}\\
x &= 5 &&\justif{on divise par $6$}
\end{calculs}
\explique{$\res{S=\ens{5}}$.}{contrôle : $30\div 5=6$}
\cas{c)}
\begin{calculs}
\dfrac{8}{x} &= -2 &&\justif{on recopie}\\
8 &= -2x &&\justif{on multiplie par $x$}\\
x &= -4 &&\justif{on divise par $-2$ : $8\div(-2)=-4$}
\end{calculs}
\explique{$\res{S=\ens{-4}}$.}{contrôle : $8\div(-4)=-2$}
\cas{d)}
\begin{calculs}
\dfrac{21}{x} &= 7 &&\justif{on recopie}\\
21 &= 7x &&\justif{on multiplie par $x$}\\
x &= 3 &&\justif{on divise par $7$}
\end{calculs}
\explique{$\res{S=\ens{3}}$.}{contrôle : $21\div 3=7$}
\cas{e)}
\explique{Soit $x$ le nombre d'amis. Chacun reçoit $120$ divisé par $x$ euros, donc
l'équation est $\dfrac{120}{x}=15$.}{on nomme l'inconnue, puis on traduit la phrase}
\begin{calculs}
\dfrac{120}{x} &= 15 &&\justif{on recopie l'équation}\\
120 &= 15x &&\justif{on multiplie par $x$}\\
x &= 8 &&\justif{on divise par $15$ : $15\times 8=120$}
\end{calculs}
\explique{La cagnotte est partagée entre \res{8~amis}.}{contrôle : $120\div 8=15$}
\end{exo}

\begin{exo}[Deux formules pour la piscine]
\cas{1.}
\explique{Sans carte : $6\times 10=\res{60}$~€.}{$10$ entrées à $6$~€}
\explique{Avec la carte : $40+4\times 10=40+40=\res{80}$~€.}{la carte, puis $10$
entrées à $4$~€}
\cas{2.}
\explique{Sans carte, on paie $6$~€ par entrée : $\res{6x}$ euros.}{on remplace $10$
par $x$ dans le calcul du 1.}
\explique{Avec la carte : $\res{40+4x}$ euros.}{$40$~€ une fois, plus $4$~€ par entrée}
\cas{3.}
\explique{Les deux formules coûtent autant quand les deux prix sont égaux : on résout
$6x=40+4x$.}{on écrit l'équation}
\begin{calculs}
6x &= 40+4x &&\justif{on recopie}\\
6x-4x &= 40 &&\justif{les $x$ à gauche}\\
2x &= 40 &&\\
x &= 20 &&\justif{on divise par $2$}
\end{calculs}
\explique{Les deux formules coûtent autant pour \res{20~entrées} : $120$~€ dans les deux
cas.}{contrôle : $6\times 20=120$ et $40+80=120$}
\cas{4.}
\explique{Sans carte : $6\times 25=150$~€. Avec la carte : $40+4\times 25=40+100=140$~€.}{on
remplace $x$ par $25$}
\explique{Pour $25$ entrées, il vaut mieux choisir \res{la carte annuelle}, moins chère
de $10$~€.}{on répond par une phrase}
\rem{Au-delà de $20$ entrées, la carte devient la formule la moins chère.}
\end{exo}

\partie{Équations $x^{2}=a$}

\begin{exo}[Résoudre $x^{2}=a$]
\explique{On regarde le signe de $a$ : si $a>0$, deux solutions $-\sqrt{a}$ et
$\sqrt{a}$ ; si $a=0$, une seule solution $0$ ; si $a<0$, aucune
solution.}{méthode 4 du cours}
\cas{a)}
\explique{$16>0$ et $16=4^{2}$ : $\res{S=\ens{-4\,;4}}$.}{$(-4)^{2}=16$ aussi}
\cas{b)}
\explique{$1>0$ et $1=1^{2}$ : $\res{S=\ens{-1\,;1}}$.}{$(-1)^{2}=1$}
\cas{c)}
\explique{Le seul nombre dont le carré vaut $0$ est $0$ : $\res{S=\ens{0}}$.}{$a=0$ : une
seule solution}
\cas{d)}
\explique{$3>0$ : $\res{S=\ens{-\sqrt{3}\,;\sqrt{3}}}$.}{$3$ n'est pas le carré d'un
entier : on garde la racine carrée}
\cas{e)}
\explique{$-25<0$ : $\res{S=\varnothing}$.}{un carré n'est jamais négatif}
\cas{f)}
\explique{$144>0$ et $144=12^{2}$ : $\res{S=\ens{-12\,;12}}$.}{$12\times 12=144$}
\rem{Erreur fréquente : oublier la solution négative. $(-4)^{2}=(-4)\times(-4)=16$.}
\rem{En e), on n'écrit pas $-5$ : $(-5)^{2}=25$, pas $-25$.}
\end{exo}

\begin{exo}[Isoler $x^{2}$ d'abord]
\explique{On isole d'abord $x^{2}$ (comme on isolerait $x$), puis on applique la méthode
de l'exercice 13.}{méthode 4 du cours}
\cas{a)}
\begin{calculs}
x^{2}-9 &= 0 &&\justif{on recopie}\\
x^{2} &= 9 &&\justif{on ajoute $9$ aux deux membres}
\end{calculs}
\explique{$9>0$ et $9=3^{2}$ : $\res{S=\ens{-3\,;3}}$.}{deux solutions opposées}
\cas{b)}
\begin{calculs}
x^{2}+4 &= 40 &&\justif{on recopie}\\
x^{2} &= 36 &&\justif{on retranche $4$}
\end{calculs}
\explique{$36=6^{2}$ : $\res{S=\ens{-6\,;6}}$.}{}
\cas{c)}
\begin{calculs}
2x^{2} &= 98 &&\justif{on recopie}\\
x^{2} &= 49 &&\justif{on divise par $2$}
\end{calculs}
\explique{$49=7^{2}$ : $\res{S=\ens{-7\,;7}}$.}{}
\cas{d)}
\begin{calculs}
5x^{2}-20 &= 0 &&\justif{on recopie}\\
5x^{2} &= 20 &&\justif{on ajoute $20$}\\
x^{2} &= 4 &&\justif{on divise par $5$}
\end{calculs}
\explique{$4=2^{2}$ : $\res{S=\ens{-2\,;2}}$.}{}
\cas{e)}
\begin{calculs}
x^{2}+10 &= 1 &&\justif{on recopie}\\
x^{2} &= -9 &&\justif{on retranche $10$ : $1-10=-9$}
\end{calculs}
\explique{$-9<0$ : $\res{S=\varnothing}$.}{un carré n'est jamais négatif}
\cas{f)}
\begin{calculs}
4x^{2}-1 &= 99 &&\justif{on recopie}\\
4x^{2} &= 100 &&\justif{on ajoute $1$}\\
x^{2} &= 25 &&\justif{on divise par $4$}
\end{calculs}
\explique{$25=5^{2}$ : $\res{S=\ens{-5\,;5}}$.}{contrôle : $4\times 25-1=99$}
\rem{Ordre des étapes : d'abord le nombre ajouté ou retranché, puis le nombre qui
multiplie $x^{2}$, comme pour une équation du premier degré.}
\end{exo}

\begin{exo}[Des carrés]
\cas{1.}
\explique{Soit $x$ le côté du potager, en mètres. L'aire d'un carré est côté $\times$
côté : $x^{2}=81$.}{on nomme l'inconnue}
\explique{$81=9^{2}$, donc $x=-9$ ou $x=9$.}{les deux solutions de l'équation}
\explique{Une longueur est positive : le côté mesure \res{9~m}.}{on rejette $-9$}
\cas{2.}
\explique{Soit $x$ le côté de la photo, en centimètres : $x^{2}=100$, donc $x=-10$ ou
$x=10$.}{$100=10^{2}$}
\explique{Le côté mesure \res{10~cm}.}{une longueur est positive}
\cas{3.}
\explique{On résout $x^{2}=49$ : $49=7^{2}$, donc les nombres cherchés sont \res{-7} et
\res{7}.}{ici, pas de longueur : on garde les deux solutions}
\rem{Dans un problème, on garde ou on rejette la solution négative selon ce que représente
$x$.}
\end{exo}

\partie{Équations produit}

\begin{exo}[Produit nul]
\explique{L'équation est de la forme \og produit $=0$\fg. Un produit est nul quand l'un
de ses facteurs est nul : on écrit \og premier facteur $=0$ ou second facteur
$=0$\fg.}{méthode 5 du cours}
\cas{a)}
\begin{calculs}
(x-1)(x+7) &= 0 &&\justif{on recopie}\\
x-1 &= 0 \quad\text{ou}\quad x+7=0 &&\justif{un des facteurs est nul}\\
x &= 1 \quad\text{ou}\quad x=-7 &&
\end{calculs}
\explique{$\res{S=\ens{-7\,;1}}$.}{solutions rangées dans l'ordre croissant}
\cas{b)}
\begin{calculs}
(x+2)(x+5) &= 0 &&\justif{on recopie}\\
x+2 &= 0 \quad\text{ou}\quad x+5=0 &&\justif{un des facteurs est nul}\\
x &= -2 \quad\text{ou}\quad x=-5 &&
\end{calculs}
\explique{$\res{S=\ens{-5\,;-2}}$.}{$-5$ est plus petit que $-2$}
\cas{c)}
\begin{calculs}
(x-3)(x-8) &= 0 &&\justif{on recopie}\\
x-3 &= 0 \quad\text{ou}\quad x-8=0 &&\\
x &= 3 \quad\text{ou}\quad x=8 &&
\end{calculs}
\explique{$\res{S=\ens{3\,;8}}$.}{}
\cas{d)}
\begin{calculs}
x(x+4) &= 0 &&\justif{les deux facteurs sont $x$ et $x+4$}\\
x &= 0 \quad\text{ou}\quad x+4=0 &&\\
x &= 0 \quad\text{ou}\quad x=-4 &&
\end{calculs}
\explique{$\res{S=\ens{-4\,;0}}$.}{}
\rem{Repère rapide : dans $(x-1)$, la solution est $1$ ; dans $(x+7)$, c'est $-7$. Le
signe change.}
\rem{Contrôle : on remplace $x$ par la solution, une parenthèse doit valoir $0$. Pour
$x=-7$ : $x+7=0$.}
\end{exo}

\begin{exo}[Des coefficients devant $x$]
\explique{Même méthode ; chaque facteur donne une petite équation du premier degré, à
résoudre jusqu'au bout.}{méthode 5 du cours}
\cas{a)}
\begin{calculs}
(2x-4)(x+3) &= 0 &&\justif{on recopie}\\
2x-4 &= 0 \quad\text{ou}\quad x+3=0 &&\justif{un des facteurs est nul}\\
2x &= 4 \quad\text{ou}\quad x=-3 &&\\
x &= 2 \quad\text{ou}\quad x=-3 &&\justif{$4\div 2=2$}
\end{calculs}
\explique{$\res{S=\ens{-3\,;2}}$.}{}
\cas{b)}
\begin{calculs}
(3x+9)(x-5) &= 0 &&\justif{on recopie}\\
3x+9 &= 0 \quad\text{ou}\quad x-5=0 &&\\
3x &= -9 \quad\text{ou}\quad x=5 &&\justif{on retranche $9$}\\
x &= -3 \quad\text{ou}\quad x=5 &&\justif{$-9\div 3=-3$}
\end{calculs}
\explique{$\res{S=\ens{-3\,;5}}$.}{}
\cas{c)}
\begin{calculs}
(5x-15)(2x+8) &= 0 &&\justif{on recopie}\\
5x-15 &= 0 \quad\text{ou}\quad 2x+8=0 &&\\
5x &= 15 \quad\text{ou}\quad 2x=-8 &&\\
x &= 3 \quad\text{ou}\quad x=-4 &&\justif{$15\div 5=3$ et $-8\div 2=-4$}
\end{calculs}
\explique{$\res{S=\ens{-4\,;3}}$.}{}
\cas{d)}
\begin{calculs}
(-x+6)(4x-8) &= 0 &&\justif{on recopie}\\
-x+6 &= 0 \quad\text{ou}\quad 4x-8=0 &&\\
6 &= x \quad\text{ou}\quad 4x=8 &&\justif{on fait passer $-x$ à droite}\\
x &= 6 \quad\text{ou}\quad x=2 &&
\end{calculs}
\explique{$\res{S=\ens{2\,;6}}$.}{}
\cas{e)}
\begin{calculs}
(-2x+10)(3x+3) &= 0 &&\justif{on recopie}\\
-2x+10 &= 0 \quad\text{ou}\quad 3x+3=0 &&\\
-2x &= -10 \quad\text{ou}\quad 3x=-3 &&\\
x &= 5 \quad\text{ou}\quad x=-1 &&\justif{$-10\div(-2)=5$ et $-3\div 3=-1$}
\end{calculs}
\explique{$\res{S=\ens{-1\,;5}}$.}{}
\rem{Avec un coefficient devant $x$, on ne lit plus la solution \og en changeant le
signe\fg{} : $2x-4=0$ donne $x=2$, pas $x=4$.}
\rem{Contrôle de e) : pour $x=5$, $-2\times 5+10=0$ ; pour $x=-1$, $3\times(-1)+3=0$.}
\end{exo}

\begin{exo}[Factoriser par $x$ d'abord]
\explique{$x^{2}-4x$ est une somme, pas un produit : on ne peut pas encore utiliser le
produit nul. Mais $x$ est dans les deux termes ($x^{2}=x\times x$) : on le met en
facteur.}{méthode 5 du cours}
\cas{a)}
\begin{calculs}
x^{2}-4x &= 0 &&\justif{on recopie ; ce n'est pas encore un produit}\\
x(x-4) &= 0 &&\justif{$x^{2}=x\times x$ et $4x=x\times 4$}\\
x &= 0 \quad\text{ou}\quad x-4=0 &&\justif{un des facteurs est nul}\\
x &= 0 \quad\text{ou}\quad x=4 &&
\end{calculs}
\explique{$\res{S=\ens{0\,;4}}$.}{}
\cas{b)}
\begin{calculs}
x^{2}+6x &= 0 &&\justif{on recopie}\\
x(x+6) &= 0 &&\justif{on factorise par $x$}\\
x &= 0 \quad\text{ou}\quad x+6=0 &&\\
x &= 0 \quad\text{ou}\quad x=-6 &&
\end{calculs}
\explique{$\res{S=\ens{-6\,;0}}$.}{}
\cas{c)}
\begin{calculs}
3x^{2}-9x &= 0 &&\justif{on recopie}\\
x(3x-9) &= 0 &&\justif{$3x^{2}=x\times 3x$ et $9x=x\times 9$}\\
x &= 0 \quad\text{ou}\quad 3x-9=0 &&\\
x &= 0 \quad\text{ou}\quad x=3 &&\justif{$3x=9$}
\end{calculs}
\explique{$\res{S=\ens{0\,;3}}$.}{}
\cas{d)}
\begin{calculs}
2x^{2}+10x &= 0 &&\justif{on recopie}\\
x(2x+10) &= 0 &&\justif{on factorise par $x$}\\
x &= 0 \quad\text{ou}\quad 2x+10=0 &&\\
x &= 0 \quad\text{ou}\quad x=-5 &&\justif{$2x=-10$}
\end{calculs}
\explique{$\res{S=\ens{-5\,;0}}$.}{}
\cas{e)}
\begin{calculs}
x^{2} &= 7x &&\justif{on recopie}\\
x^{2}-7x &= 0 &&\justif{on retranche $7x$ : tout passe à gauche}\\
x(x-7) &= 0 &&\justif{on factorise par $x$}\\
x &= 0 \quad\text{ou}\quad x-7=0 &&\\
x &= 0 \quad\text{ou}\quad x=7 &&
\end{calculs}
\explique{$\res{S=\ens{0\,;7}}$.}{contrôle : $7^{2}=49$ et $7\times 7=49$}
\rem{Contrôle de la factorisation : on redéveloppe, $x(x-4)=x^{2}-4x$.}
\rem{Piège en e) : diviser les deux membres par $x$ donne $x=7$ et fait perdre la solution
$0$. On ne divise jamais par $x$ : on factorise.}
\end{exo}

\partie{Fonction polynôme du second degré}

\begin{exo}[Reconnaître, lire $a$, $b$ et $c$]
\explique{Une fonction est du second degré si elle s'écrit $ax^{2}+bx+c$ avec $a\ne 0$ :
il faut un terme en $x^{2}$. On range les termes dans cet ordre, puis on lit chaque
coefficient \textbf{avec son signe}.}{méthode 6 du cours}
\cas{$f$}
\explique{$f(x)=2x^{2}+5x-1$ : \res{a=2} ; \res{b=5} ; \res{c=-1}.}{déjà rangée ; $c$
garde son signe $-$}
\cas{$g$}
\explique{$g(x)=-x^{2}+3$ : \res{a=-1} ; \res{b=0} ; \res{c=3}.}{$-x^{2}=-1\times x^{2}$ ;
pas de terme en $x$ : $b=0$}
\cas{$h$}
\explique{$h(x)=7x-2$ n'a pas de terme en $x^{2}$ : $h$ \res{n'est pas du second degré}
(c'est une fonction affine).}{sa courbe est une droite}
\cas{$k$}
\explique{On range : $k(x)=-x^{2}+4x$, donc \res{a=-1} ; \res{b=4} ;
\res{c=0}.}{pas de nombre seul : $c=0$}
\cas{$m$}
\explique{$m(x)=x^{2}=1x^{2}+0x+0$ : \res{a=1} ; \res{b=0} ; \res{c=0}.}{$x^{2}$, c'est
$1\times x^{2}$}
\cas{$p$}
\explique{On range : $p(x)=3x^{2}-2x+6$, donc \res{a=3} ; \res{b=-2} ;
\res{c=6}.}{chaque terme garde son signe en changeant de place}
\rem{Erreur fréquente : lire les coefficients dans l'ordre où ils sont écrits. Pour $p$,
$a$ n'est pas $6$ : $a$ est toujours le nombre devant $x^{2}$.}
\end{exo}

\begin{exo}[Calculer des images]
\explique{Pour calculer une image, on remplace $x$ par le nombre, \emph{entre
parenthèses} s'il est négatif ; on calcule d'abord les carrés, puis les produits, puis
les sommes.}{méthode 6 du cours}
\cas{1.}
\begin{calculs}
f(0) &= 0^{2}-3\times 0+2 &&\justif{on remplace $x$ par $0$}\\
f(0) &= \res{2} &&
\end{calculs}
\begin{calculs}
f(1) &= 1^{2}-3\times 1+2 &&\justif{on remplace $x$ par $1$}\\
f(1) &= 1-3+2 = \res{0} &&
\end{calculs}
\begin{calculs}
f(-1) &= (-1)^{2}-3\times(-1)+2 &&\justif{on remplace $x$ par $(-1)$}\\
f(-1) &= 1+3+2 = \res{6} &&\justif{$(-1)^{2}=1$ et $-3\times(-1)=+3$}
\end{calculs}
\begin{calculs}
f(3) &= 3^{2}-3\times 3+2 &&\justif{on remplace $x$ par $3$}\\
f(3) &= 9-9+2 = \res{2} &&
\end{calculs}
\begin{calculs}
f(-2) &= (-2)^{2}-3\times(-2)+2 &&\justif{on remplace $x$ par $(-2)$}\\
f(-2) &= 4+6+2 = \res{12} &&\justif{$(-2)^{2}=4$ et $-3\times(-2)=+6$}
\end{calculs}
\cas{2.}
\begin{calculs}
g(0) &= -2\times 0^{2}+0+5 &&\justif{on remplace $x$ par $0$}\\
g(0) &= \res{5} &&
\end{calculs}
\begin{calculs}
g(2) &= -2\times 2^{2}+2+5 &&\justif{on remplace $x$ par $2$}\\
g(2) &= -8+2+5 = \res{-1} &&\justif{$2^{2}=4$, puis $-2\times 4=-8$}
\end{calculs}
\begin{calculs}
g(-1) &= -2\times(-1)^{2}+(-1)+5 &&\justif{on remplace $x$ par $(-1)$}\\
g(-1) &= -2-1+5 = \res{2} &&\justif{$(-1)^{2}=1$, puis $-2\times 1=-2$}
\end{calculs}
\rem{Erreur fréquente : calculer $-2\times(-1)$ avant le carré, et trouver $+2$. Le carré
passe d'abord : $(-1)^{2}=1$, puis $-2\times 1=-2$.}
\rem{L'image de $0$ est toujours le nombre $c$ : $f(0)=2$ et $g(0)=5$.}
\end{exo}

\begin{exo}[Sur la courbe ou non ?]
\explique{Un point est sur la courbe de $f$ quand son ordonnée est égale à l'image de son
abscisse. On calcule donc $f(\text{abscisse})$ et on compare à
l'ordonnée.}{méthode 6 du cours}
\cas{$A(1\,;0)$}
\begin{calculs}
f(1) &= 1^{2}+2\times 1-3 &&\justif{on remplace $x$ par l'abscisse $1$}\\
f(1) &= 1+2-3 = 0 &&
\end{calculs}
\explique{$f(1)=0$, l'ordonnée de $A$ : $A$ \res{est sur} la courbe.}{}
\cas{$B(-1\,;-2)$}
\begin{calculs}
f(-1) &= (-1)^{2}+2\times(-1)-3 &&\justif{on remplace $x$ par $(-1)$}\\
f(-1) &= 1-2-3 = -4 &&
\end{calculs}
\explique{$-4\ne -2$ : $B$ \res{n'est pas} sur la courbe.}{le point de la courbe
d'abscisse $-1$ est $(-1\,;-4)$}
\cas{$C(2\,;5)$}
\begin{calculs}
f(2) &= 2^{2}+2\times 2-3 &&\\
f(2) &= 4+4-3 = 5 &&
\end{calculs}
\explique{$f(2)=5$ : $C$ \res{est sur} la courbe.}{}
\cas{$D(-3\,;0)$}
\begin{calculs}
f(-3) &= (-3)^{2}+2\times(-3)-3 &&\\
f(-3) &= 9-6-3 = 0 &&\justif{$(-3)^{2}=9$}
\end{calculs}
\explique{$f(-3)=0$ : $D$ \res{est sur} la courbe.}{}
\cas{$E(0\,;3)$}
\explique{$f(0)=-3$ et $-3\ne 3$ : $E$ \res{n'est pas} sur la courbe.}{$f(0)=c=-3$}
\rem{On ne remplace que $x$ : c'est l'abscisse (le premier nombre) qui entre dans le
calcul, l'ordonnée sert à comparer.}
\end{exo}

\begin{exo}[Distance d'arrêt]
\cas{1.}
\explique{$d(v)=v^{2}+3v$ est de la forme $av^{2}+bv+c$ avec un terme en $v^{2}$ : $d$
est du second degré, avec \res{a=1} ; \res{b=3} ; \res{c=0}.}{ici la variable s'appelle
$v$ au lieu de $x$ ; pas de nombre seul : $c=0$}
\cas{2.}
\begin{calculs}
d(3) &= 3^{2}+3\times 3 &&\justif{on remplace $v$ par $3$}\\
d(3) &= 9+9 = \res{18} &&
\end{calculs}
\begin{calculs}
d(5) &= 5^{2}+3\times 5 &&\justif{on remplace $v$ par $5$}\\
d(5) &= 25+15 = \res{40} &&
\end{calculs}
\begin{calculs}
d(9) &= 9^{2}+3\times 9 &&\justif{on remplace $v$ par $9$}\\
d(9) &= 81+27 = \res{108} &&
\end{calculs}
\explique{Interprétation : à $50$~km/h, la voiture a besoin de $40$~m pour
s'arrêter.}{$v=5$ correspond à $50$~km/h ; $d$ est en mètres}
\cas{3.}
\begin{calculs}
d(6) &= 6^{2}+3\times 6 &&\justif{$60$~km/h : $v=6$}\\
d(6) &= 36+18 = \res{54} &&
\end{calculs}
\explique{À $30$~km/h, $d(3)=18$~m ; le double serait $36$~m. Or $d(6)=54$~m : la
distance d'arrêt \res{ne double pas}, elle est multipliée par $3$.}{$54\div 18=3$}
\rem{Ce n'est pas une situation de proportionnalité : à cause du terme en $v^{2}$,
doubler la vitesse fait bien plus que doubler la distance.}
\end{exo}

\partie{Forme factorisée et racines}

\begin{exo}[Vérifier une forme factorisée]
\explique{\og Vérifier\fg{} : on part de la forme factorisée, la plus compliquée, on la
développe, et on doit retrouver l'autre expression. On développe d'abord les deux
parenthèses, puis on multiplie par le nombre devant.}{méthode 7 du cours}
\cas{a)}
\begin{calculs}
(x-2)(x+4) &= x\times x+x\times 4+(-2)\times x+(-2)\times 4 &&\justif{quatre produits}\\
(x-2)(x+4) &= x^{2}+4x-2x-8 &&\\
(x-2)(x+4) &= x^{2}+2x-8 &&\justif{$4x-2x=2x$}
\end{calculs}
\explique{On retrouve $f(x)$ : l'égalité est \res{vraie}.}{}
\cas{b)}
\begin{calculs}
3(x+2)(x-3) &= 3\left(x^{2}-3x+2x-6\right) &&\justif{d'abord les deux parenthèses}\\
3(x+2)(x-3) &= 3\left(x^{2}-x-6\right) &&\justif{$-3x+2x=-x$}\\
3(x+2)(x-3) &= 3x^{2}-3x-18 &&\justif{puis $3$ multiplie chaque terme}
\end{calculs}
\explique{On retrouve $g(x)$ : l'égalité est \res{vraie}.}{}
\cas{c)}
\begin{calculs}
-(x-1)(x-5) &= -\left(x^{2}-5x-x+5\right) &&\justif{d'abord les deux parenthèses}\\
-(x-1)(x-5) &= -\left(x^{2}-6x+5\right) &&\justif{$(-1)\times(-5)=+5$}\\
-(x-1)(x-5) &= -x^{2}+6x-5 &&\justif{le $-$ devant change tous les signes}
\end{calculs}
\explique{On retrouve $h(x)$ : l'égalité est \res{vraie}.}{}
\cas{d)}
\begin{calculs}
2(x-2)(x+2) &= 2\left(x^{2}-2^{2}\right) &&\justif{$(a-b)(a+b)=a^{2}-b^{2}$, avec $a=x$
et $b=2$}\\
2(x-2)(x+2) &= 2\left(x^{2}-4\right) &&\\
2(x-2)(x+2) &= 2x^{2}-8 &&\justif{$2$ multiplie chaque terme}
\end{calculs}
\explique{On retrouve $k(x)$ : l'égalité est \res{vraie}.}{}
\rem{Erreur fréquente : multiplier le nombre devant dans les \emph{deux} parenthèses,
$3(x+2)(x-3)=(3x+6)(3x-9)$. Le $3$ ne multiplie qu'une seule fois : on le garde devant
et on le distribue à la fin.}
\end{exo}

\begin{exo}[Lire les racines]
\explique{Les racines sont les solutions de $f(x)=0$. Sur une forme factorisée, chaque
racine annule une parenthèse ; le nombre devant les parenthèses n'est jamais
nul.}{propriété 6 du cours}
\cas{$f$}
\explique{$x-3=0$ ou $x+1=0$ : les racines de $f$ sont \res{-1} et
\res{3}.}{dans $(x+1)$, la racine est $-1$}
\cas{$g$}
\explique{$x+5=0$ ou $x-2=0$ : les racines de $g$ sont \res{-5} et \res{2}.}{le $2$
devant ne change pas les racines}
\cas{$h$}
\explique{$x-4=0$ ou $x-6=0$ : les racines de $h$ sont \res{4} et \res{6}.}{le signe $-$
devant ne change pas les racines}
\cas{$k$}
\explique{$5x=0$ ou $x-7=0$ : les racines de $k$ sont \res{0} et \res{7}.}{$5x=0$ donne
$x=0$}
\cas{$m$}
\explique{$x+2=0$ ou $x+8=0$ : les racines de $m$ sont \res{-8} et \res{-2}.}{deux
racines négatives}
\rem{Erreur fréquente : garder le signe de la parenthèse. Dans $(x+5)$, la racine est
$-5$ : on vérifie que $-5+5=0$.}
\rem{Le nombre devant les parenthèses ($2$, $-1$, $5$, $-3$) n'est pas une racine.}
\end{exo}

\begin{exo}[Points sur les axes]
\cas{1.}
\begin{calculs}
(x+2)(x-3) &= x\times x+x\times(-3)+2\times x+2\times(-3) &&\justif{quatre produits}\\
(x+2)(x-3) &= x^{2}-3x+2x-6 &&\\
(x+2)(x-3) &= x^{2}-x-6 &&\justif{$-3x+2x=-x$}
\end{calculs}
\explique{On retrouve $f(x)$ : donc \res{f(x)=(x+2)(x-3)}.}{on part de la forme
factorisée}
\cas{2.}
\begin{calculs}
(x+2)(x-3) &= 0 &&\justif{on résout $f(x)=0$ avec la forme factorisée}\\
x+2 &= 0 \quad\text{ou}\quad x-3=0 &&\justif{produit nul}\\
x &= -2 \quad\text{ou}\quad x=3 &&
\end{calculs}
\explique{Les racines de $f$ sont \res{-2} et \res{3}.}{}
\cas{3.}
\explique{La courbe coupe l'axe des abscisses aux points \res{(-2\,;0)} et
\res{(3\,;0)}.}{sur l'axe des abscisses, l'ordonnée est nulle}
\cas{4.}
\begin{calculs}
f(0) &= 0^{2}-0-6 &&\justif{la forme développée est la plus simple ici}\\
f(0) &= -6 &&
\end{calculs}
\explique{La courbe coupe l'axe des ordonnées au point \res{(0\,;-6)}.}{sur l'axe des
ordonnées, l'abscisse est nulle ; $f(0)=c$}
\rem{Chaque forme a son usage : la forme factorisée donne les racines, la forme développée
donne $f(0)=c$.}
\end{exo}

\partie{Représentation graphique}

\begin{exo}[Tableau de valeurs et tracé]
\explique{On calcule chaque image en remplaçant $x$, on place les points (en croix), puis
on les relie par une courbe arrondie, sans règle.}{méthode 8 du cours}
\cas{a)}
\begin{minipage}[t]{0.58\linewidth}\raggedright
\begin{calculs}
f(-1) &= 1+4+3 = 8 &&\justif{$(-1)^{2}=1$ et $-4\times(-1)=+4$}\\
f(0) &= 0-0+3 = 3 &&\\
f(1) &= 1-4+3 = 0 &&\\
f(2) &= 4-8+3 = -1 &&\\
f(3) &= 9-12+3 = 0 &&\\
f(4) &= 16-16+3 = 3 &&\\
f(5) &= 25-20+3 = 8 &&
\end{calculs}
\begin{center}\small
\renewcommand{\arraystretch}{1.3}
\begin{tabular}{|c|*{7}{c|}}
\hline
$x$ & $-1$ & $0$ & $1$ & $2$ & $3$ & $4$ & $5$\\
\hline
$f(x)$ & $8$ & $3$ & $0$ & $-1$ & $0$ & $3$ & $8$\\
\hline
\end{tabular}
\end{center}
\explique{$a=1>0$ : la courbe est tournée vers le haut.}{contrôle de l'allure}
\explique{Sommet \res{S(2\,;-1)} ; axe de symétrie : la droite verticale qui passe par
\res{x=2}.}{le point le plus bas ; les valeurs se répètent de part et d'autre de $2$}
\end{minipage}\hfill
\begin{minipage}[t]{0.38\linewidth}\centering\vspace{0pt}
\begin{tikzpicture}[x=0.45cm,y=0.45cm]
  \repere{-2}{-2}{6}{9}
  \gradx{-1,1,2,3,4,5} \grady{2,4,6,8}
  \draw[red,dashed,thick] (2,-2)--(2,9);
  \draw[coulCourbe,line width=1.1pt,domain=-1:5,samples=60] plot (\x,{\x*\x-4*\x+3});
  \foreach \p in {(-1,8),(0,3),(1,0),(2,-1),(3,0),(4,3),(5,8)}
    \path[coulCourbe] \p node[croix]{};
  \node[below right,font=\footnotesize,inner sep=2pt] at (2,-1) {$S$};
  \node[coulCourbe,font=\footnotesize] at (5.5,6.5) {$\Cf$};
\end{tikzpicture}
\end{minipage}
\cas{b)}
\begin{minipage}[t]{0.58\linewidth}\raggedright
\begin{calculs}
g(-2) &= -4-4+3 = -5 &&\justif{$-(-2)^{2}=-4$ et $2\times(-2)=-4$}\\
g(-1) &= -1-2+3 = 0 &&\justif{$-(-1)^{2}=-1$}\\
g(0) &= 0+0+3 = 3 &&\\
g(1) &= -1+2+3 = 4 &&\\
g(2) &= -4+4+3 = 3 &&\\
g(3) &= -9+6+3 = 0 &&\\
g(4) &= -16+8+3 = -5 &&
\end{calculs}
\begin{center}\small
\renewcommand{\arraystretch}{1.3}
\begin{tabular}{|c|*{7}{c|}}
\hline
$x$ & $-2$ & $-1$ & $0$ & $1$ & $2$ & $3$ & $4$\\
\hline
$g(x)$ & $-5$ & $0$ & $3$ & $4$ & $3$ & $0$ & $-5$\\
\hline
\end{tabular}
\end{center}
\explique{$a=-1<0$ : la courbe est tournée vers le bas.}{contrôle de l'allure}
\explique{Sommet \res{S(1\,;4)} ; axe de symétrie : la droite verticale qui passe par
\res{x=1}.}{le point le plus haut}
\end{minipage}\hfill
\begin{minipage}[t]{0.38\linewidth}\centering\vspace{0pt}
\begin{tikzpicture}[x=0.45cm,y=0.45cm]
  \repere{-3}{-6}{5}{5}
  \gradx{-2,-1,1,2,3,4} \grady{-4,-2,2,4}
  \draw[red,dashed,thick] (1,-6)--(1,5);
  \draw[coulCourbe,line width=1.1pt,domain=-2:4,samples=60] plot (\x,{-\x*\x+2*\x+3});
  \foreach \p in {(-2,-5),(-1,0),(0,3),(1,4),(2,3),(3,0),(4,-5)}
    \path[coulCourbe] \p node[croix]{};
  \node[above right,font=\footnotesize,inner sep=2pt] at (1,4) {$S$};
  \node[coulCourbe,font=\footnotesize] at (4.3,-3) {$\mathcal{C}_{g}$};
\end{tikzpicture}
\end{minipage}
\rem{Erreur fréquente : $-x^{2}$ pour $x=-2$ vaut $-(-2)^{2}=-4$, pas $+4$. On fait le
carré d'abord, puis on met le signe $-$.}
\rem{Les valeurs égales de part et d'autre du sommet ($f(1)=f(3)$, $g(0)=g(2)$) servent
de contrôle : si la symétrie n'y est pas, un calcul est faux.}
\end{exo}

\begin{exo}[Associer une fonction à sa courbe]
\begin{minipage}[t]{0.58\linewidth}\raggedright
\explique{Premier indice, le signe de $a$ : $g$ et $k$ ont $a=-1<0$, leurs courbes sont
tournées vers le bas ($\mathcal{C}_{1}$ et $\mathcal{C}_{4}$). $f$ et $h$ ont $a>0$,
leurs courbes sont tournées vers le haut ($\mathcal{C}_{2}$ et $\mathcal{C}_{3}$).}{la
propriété 7 du cours}
\explique{Second indice, le point $(0\,;c)$ sur l'axe des ordonnées, car
$f(0)=c$.}{méthode 9 du cours}
\explique{$g(0)=4$ : $g$ correspond à $\res{\mathcal{C}_{1}}$, qui passe par
$(0\,;4)$.}{}
\explique{$k(0)=-1$ : $k$ correspond à $\res{\mathcal{C}_{4}}$, qui passe par
$(0\,;-1)$.}{}
\explique{$f(0)=-3$ : $f$ correspond à $\res{\mathcal{C}_{3}}$, qui passe par
$(0\,;-3)$.}{}
\explique{$h(0)=0$ : $h$ correspond à $\res{\mathcal{C}_{2}}$, qui passe par
l'origine.}{}
\end{minipage}\hfill
\begin{minipage}[t]{0.38\linewidth}\centering\vspace{0pt}
\begin{tikzpicture}[x=0.45cm,y=0.45cm]
  \repere{-4}{-5}{4}{5}
  \gradx{-3,-2,-1,1,2,3} \grady{-4,-2,2,4}
  \draw[coulCourbe,line width=1.1pt,domain=-3:3,samples=60] plot (\x,{-\x*\x+4});
  \draw[coulCourbeB,line width=1.1pt,domain=-1.58:1.58,samples=60] plot (\x,{2*\x*\x});
  \draw[coulCourbeC,line width=1.1pt,domain=-2.8:2.8,samples=60] plot (\x,{\x*\x-3});
  \draw[black!70,line width=1.1pt,domain=-2:2,samples=60] plot (\x,{-\x*\x-1});
  \foreach \y in {4,0,-3,-1} \path (0,\y) node[croix]{};
  \node[coulCourbe,font=\footnotesize,fill=white,inner sep=1pt] at (-3.4,-4.3) {$\mathcal{C}_{1}$};
  \node[coulCourbeB,font=\footnotesize,fill=white,inner sep=1pt] at (2.1,4.6) {$\mathcal{C}_{2}$};
  \node[coulCourbeC,font=\footnotesize,fill=white,inner sep=1pt] at (3.4,4.3) {$\mathcal{C}_{3}$};
  \node[black!70,font=\footnotesize,fill=white,inner sep=1pt] at (-2.6,-3.3) {$\mathcal{C}_{4}$};
\end{tikzpicture}
\end{minipage}
\rem{Deux indices suffisent : le signe de $a$ partage les courbes en deux groupes, puis
le point $(0\,;c)$ les distingue dans chaque groupe.}
\end{exo}

\begin{exo}[Lire une parabole]
\begin{minipage}[t]{0.58\linewidth}\raggedright
\cas{1.}
\explique{La parabole est tournée vers le \res{haut}. C'est cohérent : $f(x)=x^{2}+2x-3$
a pour coefficient $a=1>0$.}{$a>0$ : vers le haut}
\cas{2.}
\explique{Le sommet est le point le plus bas : \res{S(-1\,;-4)}.}{contrôle :
$f(-1)=1-2-3=-4$}
\explique{L'axe de symétrie est la droite verticale qui passe par
\res{x=-1}.}{elle passe par le sommet}
\cas{3.}
\explique{La courbe coupe l'axe des abscisses en $x=-3$ et $x=1$ : les racines de $f$
sont \res{-3} et \res{1}.}{contrôle : $f(1)=0$ et $f(-3)=9-6-3=0$}
\cas{4.}
\explique{On trace la droite horizontale d'ordonnée $-3$ ; elle coupe la courbe en deux
points, d'abscisses $-2$ et $0$ : $\res{S=\ens{-2\,;0}}$.}{les solutions sont des
abscisses}
\cas{5.}
\begin{calculs}
f(0) &= 0^{2}+2\times 0-3 &&\justif{on remplace $x$ par $0$}\\
f(0) &= \res{-3} &&
\end{calculs}
\explique{C'est le point $(0\,;-3)$, où la courbe coupe l'axe des ordonnées ; c'est l'un
des deux points de la question 4.}{$f(0)=c$}
\end{minipage}\hfill
\begin{minipage}[t]{0.38\linewidth}\centering\vspace{0pt}
\begin{tikzpicture}[x=0.45cm,y=0.45cm]
  \repere{-5}{-5}{3}{6}
  \gradx{-4,-3,-2,-1,1,2} \grady{-4,-3,-2,-1,1,2,3,4,5}
  \draw[red,dashed,thick] (-1,-5)--(-1,6);
  \draw[coulExo,thick] (-5,-3)--(3,-3);
  \draw[coulCourbe,line width=1.1pt,domain=-4:2,samples=60] plot (\x,{\x*\x+2*\x-3});
  \draw[densely dotted,thick] (-2,-3)--(-2,0);
  \foreach \p in {(-1,-4),(-3,0),(1,0),(-2,-3),(0,-3)}
    \path \p node[croix]{};
  \node[below right,font=\footnotesize,inner sep=2pt] at (-1,-4) {$S$};
  \node[coulExo,font=\footnotesize,above] at (-4.2,-3) {$y=-3$};
  \node[coulCourbe,font=\footnotesize] at (2.4,4) {$\Cf$};
\end{tikzpicture}
\end{minipage}
\rem{Résoudre graphiquement $f(x)=-3$ : on cherche les points de la courbe
d'\emph{ordonnée} $-3$, et on donne leurs \emph{abscisses}.}
\end{exo}

\partie{Tableaux de variations}

\begin{exo}[Sommet et tableau de variations]
\explique{On lit les racines sur la forme factorisée ; le sommet est au milieu des
racines : $x_{S}=\dfrac{x_{1}+x_{2}}{2}$, et son ordonnée est $f(x_{S})$. Le signe de
$a$, le nombre devant les parenthèses, donne le sens du tableau.}{méthode 10 du cours}
\cas{a)}
\explique{Racines de $f$ : $2$ et $6$.}{on annule chaque parenthèse}
\begin{calculs}
x_{S} &= \dfrac{2+6}{2} &&\justif{le milieu des racines}\\
x_{S} &= 4 &&\\
f(4) &= (4-2)\times(4-6) &&\justif{on remplace $x$ par $4$ dans la forme factorisée}\\
f(4) &= 2\times(-2) = -4 &&
\end{calculs}
\explique{Sommet \res{S(4\,;-4)}.}{}
\explique{Pas de nombre devant les parenthèses : $a=1>0$. $f$ diminue, puis augmente ;
son minimum est $-4$.}{tournée vers le haut}
\varmin{f(x)}{4}{-4}
\cas{b)}
\explique{Racines de $g$ : $-1$ et $3$.}{dans $(x+1)$, la racine est $-1$}
\begin{calculs}
x_{S} &= \dfrac{-1+3}{2} &&\justif{le milieu des racines}\\
x_{S} &= 1 &&\\
g(1) &= -2\times(1+1)\times(1-3) &&\justif{on remplace $x$ par $1$}\\
g(1) &= -2\times 2\times(-2) = 8 &&\justif{deux signes $-$ : le produit est positif}
\end{calculs}
\explique{Sommet \res{S(1\,;8)}.}{}
\explique{$a=-2<0$ : $g$ augmente, puis diminue ; son maximum est $8$.}{tournée vers le
bas}
\varmax{g(x)}{1}{8}
\cas{c)}
\explique{Racines de $h$ : $-4$ et $-2$.}{deux racines négatives}
\begin{calculs}
x_{S} &= \dfrac{-4+(-2)}{2} &&\justif{le milieu des racines}\\
x_{S} &= \dfrac{-6}{2} = -3 &&\\
h(-3) &= 3\times(-3+4)\times(-3+2) &&\justif{on remplace $x$ par $(-3)$}\\
h(-3) &= 3\times 1\times(-1) = -3 &&
\end{calculs}
\explique{Sommet \res{S(-3\,;-3)}.}{}
\explique{$a=3>0$ : $h$ diminue, puis augmente ; son minimum est $-3$.}{tournée vers le
haut}
\varmin{h(x)}{-3}{-3}
\cas{d)}
\explique{Racines de $k$ : $5$ et $-1$.}{}
\begin{calculs}
x_{S} &= \dfrac{5+(-1)}{2} &&\justif{le milieu des racines}\\
x_{S} &= \dfrac{4}{2} = 2 &&\\
k(2) &= -(2-5)\times(2+1) &&\justif{on remplace $x$ par $2$}\\
k(2) &= -(-3)\times 3 = 9 &&\justif{$-(-3)=3$, puis $3\times 3=9$}
\end{calculs}
\explique{Sommet \res{S(2\,;9)}.}{}
\explique{$a=-1<0$ : $k$ augmente, puis diminue ; son maximum est $9$.}{le signe $-$
devant les parenthèses, c'est $a=-1$}
\varmax{k(x)}{2}{9}
\rem{$a$ est le nombre devant les parenthèses : dans $-(x-5)(x+1)$, $a=-1$ ; sans
nombre devant, $a=1$.}
\rem{Contrôle du sens : avec $a>0$, la valeur du sommet est la plus petite (souvent
négative) ; avec $a<0$, la plus grande.}
\end{exo}

\begin{exo}[Factoriser d'abord]
\explique{Les fonctions sont données sous forme développée : pour trouver les racines, on
résout $f(x)=0$ en factorisant par $x$ (ou, en d), avec la méthode
$x^{2}=a$).}{méthodes 5 et 10 du cours}
\cas{a)}
\begin{calculs}
x^{2}-6x &= 0 &&\justif{les racines : on résout $f(x)=0$}\\
x(x-6) &= 0 &&\justif{on factorise par $x$}\\
x &= 0 \quad\text{ou}\quad x=6 &&\justif{produit nul}
\end{calculs}
\begin{calculs}
x_{S} &= \dfrac{0+6}{2} = 3 &&\justif{le milieu des racines}\\
f(3) &= 3^{2}-6\times 3 &&\justif{la forme développée suffit}\\
f(3) &= 9-18 = -9 &&
\end{calculs}
\explique{Sommet \res{S(3\,;-9)} ; $a=1>0$ : minimum $-9$.}{tournée vers le haut}
\varmin{f(x)}{3}{-9}
\cas{b)}
\begin{calculs}
-x^{2}+10x &= 0 &&\justif{les racines}\\
x(-x+10) &= 0 &&\justif{on factorise par $x$ : $-x^{2}=x\times(-x)$}\\
x &= 0 \quad\text{ou}\quad -x+10=0 &&\\
x &= 0 \quad\text{ou}\quad x=10 &&
\end{calculs}
\begin{calculs}
x_{S} &= \dfrac{0+10}{2} = 5 &&\justif{le milieu des racines}\\
g(5) &= -5^{2}+10\times 5 &&\justif{$-5^{2}=-25$ : le carré, puis le signe}\\
g(5) &= -25+50 = 25 &&
\end{calculs}
\explique{Sommet \res{S(5\,;25)} ; $a=-1<0$ : maximum $25$.}{tournée vers le bas}
\varmax{g(x)}{5}{25}
\cas{c)}
\begin{calculs}
2x^{2}+8x &= 0 &&\justif{les racines}\\
x(2x+8) &= 0 &&\justif{on factorise par $x$}\\
x &= 0 \quad\text{ou}\quad 2x+8=0 &&\\
x &= 0 \quad\text{ou}\quad x=-4 &&\justif{$2x=-8$}
\end{calculs}
\begin{calculs}
x_{S} &= \dfrac{-4+0}{2} = -2 &&\justif{le milieu des racines}\\
h(-2) &= 2\times(-2)^{2}+8\times(-2) &&\justif{on remplace $x$ par $(-2)$}\\
h(-2) &= 8-16 = -8 &&\justif{$(-2)^{2}=4$, puis $2\times 4=8$}
\end{calculs}
\explique{Sommet \res{S(-2\,;-8)} ; $a=2>0$ : minimum $-8$.}{tournée vers le haut}
\varmin{h(x)}{-2}{-8}
\cas{d)}
\begin{calculs}
x^{2}-9 &= 0 &&\justif{les racines ; pas de $x$ en facteur ici}\\
x^{2} &= 9 &&\justif{on isole $x^{2}$}\\
x &= -3 \quad\text{ou}\quad x=3 &&\justif{$9>0$ et $9=3^{2}$}
\end{calculs}
\begin{calculs}
x_{S} &= \dfrac{-3+3}{2} = 0 &&\justif{le milieu des racines}\\
k(0) &= 0^{2}-9 = -9 &&
\end{calculs}
\explique{Sommet \res{S(0\,;-9)} ; $a=1>0$ : minimum $-9$.}{tournée vers le haut}
\varmin{k(x)}{0}{-9}
\rem{On peut aussi factoriser $k(x)$ avec l'identité $a^{2}-b^{2}$ :
$x^{2}-9=(x-3)(x+3)$, et lire les racines $-3$ et $3$.}
\rem{Contrôle de la factorisation par $x$ : on redéveloppe, $x(-x+10)=-x^{2}+10x$.}
\end{exo}

\begin{exo}[Maximum ou minimum ?]
\cas{1.}
\explique{Pour $f(x)=4(x-1)(x-3)$ : racines $1$ et $3$, donc
$x_{S}=\dfrac{1+3}{2}=2$.}{le milieu des racines}
\begin{calculs}
f(2) &= 4\times(2-1)\times(2-3) &&\justif{on remplace $x$ par $2$}\\
f(2) &= 4\times 1\times(-1) = -4 &&
\end{calculs}
\explique{$a=4>0$ : $f$ admet un \res{minimum, égal à $-4$}, atteint en
$x=2$.}{tournée vers le haut}
\explique{Pour $g(x)=-(x+3)(x-1)$ : racines $-3$ et $1$, donc
$x_{S}=\dfrac{-3+1}{2}=-1$.}{}
\begin{calculs}
g(-1) &= -(-1+3)\times(-1-1) &&\justif{on remplace $x$ par $(-1)$}\\
g(-1) &= -2\times(-2) = 4 &&
\end{calculs}
\explique{$a=-1<0$ : $g$ admet un \res{maximum, égal à $4$}, atteint en
$x=-1$.}{tournée vers le bas}
\explique{Pour $h(x)=-5x(x-2)$ : racines $0$ et $2$, donc
$x_{S}=\dfrac{0+2}{2}=1$.}{$-5x=0$ donne $x=0$}
\begin{calculs}
h(1) &= -5\times 1\times(1-2) &&\justif{on remplace $x$ par $1$}\\
h(1) &= -5\times(-1) = 5 &&
\end{calculs}
\explique{$a=-5<0$ : $h$ admet un \res{maximum, égal à $5$}, atteint en
$x=1$.}{tournée vers le bas}
\cas{2.}
\begin{calculs}
f(0) &= 4\times(0-1)\times(0-3) &&\justif{on remplace $x$ par $0$}\\
f(0) &= 4\times(-1)\times(-3) = \res{12} &&\justif{deux signes $-$ : positif}
\end{calculs}
\begin{calculs}
f(4) &= 4\times(4-1)\times(4-3) &&\justif{on remplace $x$ par $4$}\\
f(4) &= 4\times 3\times 1 = \res{12} &&
\end{calculs}
\explique{Entre $0$ et $4$, $f$ diminue de $12$ à $-4$ (en $x=2$), puis augmente
jusqu'à $12$.}{les valeurs aux bouts, puis le sommet}
\varint{f(x)}{$0$,$2$,$4$}{+/$12$,-/$-4$,+/$12$}
\rem{$f(0)=f(4)$ : c'est normal, $0$ et $4$ sont à la même distance du sommet $x=2$
(symétrie de la parabole).}
\end{exo}

\partie{Problèmes type épreuve anticipée}

\begin{exo}[Un food truck]
\explique{On garde en tête les unités : $x$ compte des \emph{dizaines} de repas, $B(x)$
des \emph{dizaines} d'euros.}{méthode 11 du cours}
\cas{1.}
\begin{calculs}
B(0) &= -0^{2}+12\times 0-20 &&\justif{on remplace $x$ par $0$}\\
B(0) &= \res{-20} &&
\end{calculs}
\explique{Si le food truck ne vend aucun repas, il perd $20$ dizaines d'euros, soit
$200$~€ : ce sont les frais de la journée.}{$B(0)<0$ : une perte}
\cas{2.}
\begin{calculs}
-(x-2)(x-10) &= -\left(x^{2}-10x-2x+20\right) &&\justif{d'abord les deux parenthèses}\\
-(x-2)(x-10) &= -\left(x^{2}-12x+20\right) &&\justif{$(-2)\times(-10)=+20$}\\
-(x-2)(x-10) &= -x^{2}+12x-20 &&\justif{le $-$ devant change tous les signes}
\end{calculs}
\explique{On retrouve $B(x)$ : donc $B(x)=-(x-2)(x-10)$.}{on part de la forme donnée}
\cas{3.}
\begin{calculs}
-(x-2)(x-10) &= 0 &&\justif{bénéfice nul : on résout $B(x)=0$}\\
x-2 &= 0 \quad\text{ou}\quad x-10=0 &&\justif{produit nul}\\
x &= 2 \quad\text{ou}\quad x=10 &&
\end{calculs}
\explique{Le bénéfice est nul pour \res{20} ou \res{100} repas.}{$x=2$ : $2$ dizaines de
repas}
\cas{4.}
\begin{calculs}
x_{S} &= \dfrac{2+10}{2} = 6 &&\justif{le milieu des racines}\\
B(6) &= -(6-2)\times(6-10) &&\justif{forme factorisée}\\
B(6) &= -4\times(-4) = 16 &&\\
B(12) &= -(12-2)\times(12-10) &&\justif{la valeur au bout de l'intervalle}\\
B(12) &= -10\times 2 = -20 &&
\end{calculs}
\explique{$a=-1<0$ : $B$ augmente de $0$ à $6$, puis diminue de $6$ à $12$.}{tournée
vers le bas ; $B(0)=-20$ (question 1)}
\varint{B(x)}{$0$,$6$,$12$}{-/$-20$,+/$16$,-/$-20$}
\cas{5.}
\explique{Le bénéfice maximal est de $16$ dizaines d'euros, soit \res{160~€}, pour $6$
dizaines de repas, soit \res{60~repas}.}{on revient aux unités de l'énoncé}
\rem{Erreur fréquente : répondre \og $16$~€ pour $6$ repas\fg. On relit l'énoncé pour
convertir les dizaines.}
\end{exo}

\begin{exo}[Une balle lancée d'une falaise]
\cas{1.}
\begin{calculs}
h(0) &= -5\times 0^{2}+10\times 0+15 &&\justif{on remplace $t$ par $0$}\\
h(0) &= \res{15} &&
\end{calculs}
\explique{Au départ ($t=0$), la balle est à $15$~m au-dessus de la mer : c'est la hauteur
de la falaise.}{$h(0)$ : la hauteur au départ}
\cas{2.}
\begin{calculs}
-5(t+1)(t-3) &= -5\left(t^{2}-3t+t-3\right) &&\justif{d'abord les deux parenthèses}\\
-5(t+1)(t-3) &= -5\left(t^{2}-2t-3\right) &&\justif{$-3t+t=-2t$}\\
-5(t+1)(t-3) &= -5t^{2}+10t+15 &&\justif{$-5$ multiplie chaque terme}
\end{calculs}
\explique{On retrouve $h(t)$ : donc $h(t)=-5(t+1)(t-3)$.}{$(-5)\times(-3)=+15$}
\cas{3.}
\begin{calculs}
-5(t+1)(t-3) &= 0 &&\justif{dans l'eau, la hauteur est nulle : $h(t)=0$}\\
t+1 &= 0 \quad\text{ou}\quad t-3=0 &&\justif{produit nul ; $-5$ n'est jamais nul}\\
t &= -1 \quad\text{ou}\quad t=3 &&
\end{calculs}
\explique{Un temps est positif : la balle touche l'eau au bout de \res{3~secondes}.}{on
rejette $t=-1$}
\cas{4.}
\begin{calculs}
t_{S} &= \dfrac{-1+3}{2} = 1 &&\justif{le milieu des racines, même avec $-1$}\\
h(1) &= -5\times(1+1)\times(1-3) &&\justif{forme factorisée}\\
h(1) &= -5\times 2\times(-2) = 20 &&
\end{calculs}
\explique{$h(0)=15$ et $h(3)=0$ ; $a=-5<0$ : $h$ augmente de $0$ à $1$, puis diminue de
$1$ à $3$.}{tournée vers le bas}
\varint[t]{h(t)}{$0$,$1$,$3$}{-/$15$,+/$20$,-/$0$}
\cas{5.}
\explique{La balle atteint sa hauteur maximale, \res{20~m}, au bout de \res{1~seconde}.}{le
maximum du tableau}
\rem{La racine $-1$ n'a pas de sens pour la balle (un temps négatif), mais elle sert
quand même à trouver le sommet : le milieu de $-1$ et $3$.}
\end{exo}

\begin{exo}[Le plus grand enclos]
\cas{1.}
\explique{Le périmètre, $40$~m, fait deux longueurs et deux largeurs : deux côtés voisins
font la moitié, $20$~m.}{$40\div 2=20$}
\explique{Si $x=5$, l'autre côté mesure $20-5=\res{15~m}$, et l'aire vaut
$5\times 15=\res{75~m^{2}}$.}{aire d'un rectangle : longueur $\times$ largeur}
\cas{2.}
\explique{Comme au 1., deux côtés voisins font $20$~m : si l'un mesure $x$, l'autre
mesure $20-x$.}{on remplace $5$ par $x$}
\explique{L'aire est le produit des deux côtés : $A(x)=x(20-x)$.}{}
\cas{3.}
\begin{calculs}
A(x) &= x(20-x) &&\justif{on recopie}\\
A(x) &= x\times 20+x\times(-x) &&\justif{$x$ multiplie chaque terme}\\
A(x) &= 20x-x^{2} &&\\
A(x) &= \res{-x^{2}+20x} &&\justif{on range}
\end{calculs}
\begin{calculs}
x(20-x) &= 0 &&\justif{les racines : la forme factorisée sert ici}\\
x &= 0 \quad\text{ou}\quad 20-x=0 &&\justif{produit nul}\\
x &= 0 \quad\text{ou}\quad x=20 &&
\end{calculs}
\explique{Les racines de $A$ sont \res{0} et \res{20}.}{un côté nul : l'enclos est
aplati, son aire est nulle}
\cas{4.}
\begin{calculs}
x_{S} &= \dfrac{0+20}{2} = 10 &&\justif{le milieu des racines}\\
A(10) &= 10\times(20-10) &&\\
A(10) &= 10\times 10 = 100 &&
\end{calculs}
\explique{$a=-1<0$ : $A$ augmente de $0$ à $10$, puis diminue de $10$ à $20$.}{tournée
vers le bas ; $A(0)=A(20)=0$}
\varint{A(x)}{$0$,$10$,$20$}{-/$0$,+/$100$,-/$0$}
\cas{5.}
\explique{L'aire est la plus grande pour $x=10$ : l'autre côté mesure aussi $20-10=10$~m.
C'est un \res{carré de $10$~m de côté}, d'aire \res{100~m^{2}}.}{on donne les deux
dimensions}
\end{exo}

\partiesansnum{Exercices type épreuve anticipée}

\begin{exo}[Automatismes (QCM)]
\cas{1.}
\begin{calculs}
6x+1 &= 2x+9 &&\justif{on recopie}\\
6x-2x &= 9-1 &&\justif{les $x$ à gauche, les nombres à droite}\\
4x &= 8 &&\\
x &= 2 &&\justif{on divise par $4$}
\end{calculs}
\explique{Réponse \res{b)}.}{contrôle : $6\times 2+1=13$ et $2\times 2+9=13$}
\cas{2.}
\begin{calculs}
(x+3)^{2} &= x^{2}+2\times x\times 3+3^{2} &&\justif{$(a+b)^{2}$, avec $a=x$ et $b=3$}\\
(x+3)^{2} &= x^{2}+6x+9 &&
\end{calculs}
\explique{Réponse \res{b)}.}{a) oublie le double produit}
\cas{3.}
\explique{$7>0$ : deux solutions, $-\sqrt{7}$ et $\sqrt{7}$. Réponse \res{b)}.}{c)
oublie la solution négative ; d) divise $7$ par $2$}
\cas{4.}
\begin{calculs}
(3x-6)(x+4) &= 0 &&\justif{on recopie}\\
3x-6 &= 0 \quad\text{ou}\quad x+4=0 &&\justif{produit nul}\\
x &= 2 \quad\text{ou}\quad x=-4 &&\justif{$3x=6$}
\end{calculs}
\explique{$S=\ens{-4\,;2}$ : réponse \res{b)}.}{d) oublie de diviser par $3$}
\cas{5.}
\begin{calculs}
y &= 2^{2}-3\times 2+1 &&\justif{on remplace $x$ par $2$}\\
y &= 4-6+1 = -1 &&
\end{calculs}
\explique{Le point est $(2\,;-1)$ : réponse \res{a)}.}{}
\cas{6.}
\explique{$a=-2<0$ : la parabole est tournée vers le bas, donc $f$ admet un
\emph{maximum}.}{on élimine a) et d)}
\explique{Racines $1$ et $5$ : le sommet est en $x_{S}=\dfrac{1+5}{2}=3$. Réponse
\res{b)}.}{le milieu des racines}
\end{exo}

\begin{exo}[Une épidémie de grippe]
\begin{minipage}[t]{0.5\linewidth}\raggedright
\cas{1.}
\explique{Sur la courbe, le point d'abscisse $2$ a pour ordonnée $12$ : $f(2)=\res{12}$.}{on
monte de $x=2$ jusqu'à la courbe}
\explique{Deux semaines après le début de l'épidémie, il y a $12$ centaines de malades,
soit \res{1\,200~malades}.}{$f(x)$ compte des centaines}
\cas{2.}
\explique{Le point le plus haut de la courbe est $(4\,;16)$ : le pic a lieu à la
semaine \res{4}, avec $16$ centaines, soit \res{1\,600~malades}.}{le sommet de la
parabole}
\end{minipage}\hfill
\begin{minipage}[t]{0.46\linewidth}\centering\vspace{0pt}
\begin{tikzpicture}[x=0.6cm,y=0.2cm]
  \draw[grille,xstep=1,ystep=2] (0,0) grid (9,18);
  \draw[-{Stealth[length=2mm]},thick] (0,0)--(9.4,0) node[below left=-1pt]{$x$};
  \draw[-{Stealth[length=2mm]},thick] (0,0)--(0,19.2) node[below left=-1pt]{$y$};
  \node[below left,font=\scriptsize,inner sep=1.5pt] at (0,0) {$0$};
  \foreach \i in {1,...,8} {\draw (\i,0.3)--(\i,-0.3) node[below,font=\scriptsize,inner sep=1.5pt]{$\i$};}
  \foreach \v in {2,4,...,16} {\draw (0.07,\v)--(-0.07,\v) node[left,font=\scriptsize,inner sep=1.5pt]{$\v$};}
  \draw[coulCourbe,line width=1.1pt,domain=0:8,samples=60] plot (\x,{-\x*\x+8*\x});
  \draw[densely dotted,thick] (2,0)--(2,12)--(0,12);
  \draw[densely dotted,thick] (4,0)--(4,16)--(0,16);
  \path (2,12) node[croix]{} (4,16) node[croix]{};
  \node[coulCourbe,font=\footnotesize] at (7.6,11) {$\Cf$};
\end{tikzpicture}
\end{minipage}
\cas{3.}
\begin{calculs}
f(1) &= -1^{2}+8\times 1 &&\justif{on remplace $x$ par $1$ ; $-1^{2}=-1$}\\
f(1) &= -1+8 = \res{7} &&
\end{calculs}
\explique{Une semaine après le début de l'épidémie, il y a \res{700~malades}.}{$7$
centaines}
\cas{4.}
\begin{calculs}
-x(x-8) &= (-x)\times x+(-x)\times(-8) &&\justif{$-x$ multiplie chaque terme}\\
-x(x-8) &= -x^{2}+8x &&\justif{$(-x)\times(-8)=+8x$}
\end{calculs}
\explique{On retrouve $f(x)$ : donc $f(x)=-x(x-8)$.}{}
\begin{calculs}
-x(x-8) &= 0 &&\justif{plus de malades : $f(x)=0$}\\
-x &= 0 \quad\text{ou}\quad x-8=0 &&\justif{produit nul}\\
x &= 0 \quad\text{ou}\quad x=8 &&
\end{calculs}
\explique{$x=0$ est le début de l'épidémie : elle s'arrête au bout de
\res{8~semaines}.}{on interprète chaque solution}
\cas{5.}
\begin{calculs}
x_{S} &= \dfrac{0+8}{2} = 4 &&\justif{le milieu des racines}\\
f(4) &= -4\times(4-8) &&\justif{forme factorisée}\\
f(4) &= -4\times(-4) = 16 &&
\end{calculs}
\explique{$a=-1<0$ : $f$ augmente de $0$ à $4$, puis diminue de $4$ à
$8$.}{$f(0)=f(8)=0$}
\varint{f(x)}{$0$,$4$,$8$}{-/$0$,+/$16$,-/$0$}
\explique{Au pic, il y a $16$ centaines de malades, soit \res{1\,600~malades} : on
retrouve la lecture de la question 2.}{le calcul confirme le graphique}
\end{exo}

\begin{exo}[Le prix d'une place de théâtre]
\cas{1.}
\begin{calculs}
200-10\times 8 &= 200-80 &&\justif{on remplace $x$ par $8$}\\
200-10\times 8 &= 120 &&
\end{calculs}
\explique{La troupe prévoit \res{120~spectateurs}. La recette est
$8\times 120=\res{960~€}$.}{recette : prix d'une place $\times$ nombre de spectateurs}
\cas{2.}
\begin{calculs}
200-10x &= 150 &&\justif{nombre de spectateurs égal à $150$}\\
-10x &= 150-200 &&\justif{on retranche $200$}\\
-10x &= -50 &&\\
x &= 5 &&\justif{on divise par $-10$}
\end{calculs}
\explique{Une place à \res{5~€} donne $150$ spectateurs.}{contrôle : $200-50=150$}
\cas{3.}
\explique{Recette $=$ prix d'une place $\times$ nombre de spectateurs, donc
$R(x)=x(200-10x)$.}{comme au 1., avec $x$ au lieu de $8$}
\begin{calculs}
R(x) &= x(200-10x) &&\justif{on recopie}\\
R(x) &= x\times 200+x\times(-10x) &&\justif{$x$ multiplie chaque terme}\\
R(x) &= 200x-10x^{2} &&\\
R(x) &= -10x^{2}+200x &&\justif{on range}
\end{calculs}
\cas{4.}
\begin{calculs}
x(200-10x) &= 0 &&\justif{on résout $R(x)=0$ avec la forme factorisée}\\
x &= 0 \quad\text{ou}\quad 200-10x=0 &&\justif{produit nul}\\
x &= 0 \quad\text{ou}\quad x=20 &&\justif{$10x=200$}
\end{calculs}
\explique{Les racines de $R$ sont \res{0} et \res{20}.}{}
\explique{Pour $x=0$, les places sont gratuites : la recette est nulle. Pour $x=20$, la
troupe ne prévoit plus aucun spectateur ($200-200=0$) : la recette est nulle
aussi.}{on interprète chaque racine}
\cas{5.}
\begin{calculs}
x_{S} &= \dfrac{0+20}{2} = 10 &&\justif{le milieu des racines}\\
R(10) &= 10\times(200-10\times 10) &&\justif{forme factorisée}\\
R(10) &= 10\times 100 = 1\,000 &&
\end{calculs}
\explique{$a=-10<0$ : $R$ augmente de $0$ à $10$, puis diminue de $10$ à
$20$.}{tournée vers le bas}
\varint{R(x)}{$0$,$10$,$20$}{-/$0$,+/$1\,000$,-/$0$}
\cas{6.}
\explique{La recette est la plus grande pour une place à \res{10~€} : elle vaut alors
\res{1\,000~€} ($100$ spectateurs).}{le maximum du tableau}
\rem{Un prix trop bas rapporte peu par place ; un prix trop haut fait fuir les
spectateurs : le meilleur prix est entre les deux.}
\end{exo}

\end{document}
